\documentclass[10pt,a4paper,fleqn]{article}

\usepackage[
a4paper,
top=18mm,
bottom=18mm,
left=19mm,
right=19mm
]{geometry}


\usepackage{graphicx}
\usepackage{amsmath}
\usepackage{newtxtext,newtxmath}
\usepackage{booktabs}
\usepackage{array}
\usepackage{multirow}
\usepackage{tabularx}
\usepackage{float}
\usepackage{placeins}

\usepackage[
font=footnotesize,
labelfont=bf,
skip=2pt
]{caption}
\usepackage{subcaption}

\usepackage{url}
\usepackage{hyperref}
\usepackage{authblk}
\usepackage{setspace}
\usepackage{lineno}
\usepackage{enumitem}
\usepackage{makecell}
\usepackage[round,authoryear]{natbib}
\usepackage{algorithm}
\usepackage{algpseudocode}
\usepackage{titlesec}

\algrenewcommand\algorithmicindent{1.3em}
\algrenewcommand\algorithmicrequire{\textbf{Input:}}
\algrenewcommand\algorithmicensure{\textbf{Output:}}

\urlstyle{same}
\setstretch{1.03}
\linenumbers

\setlength{\parindent}{1.2em}
\setlength{\parskip}{0pt}

\setlength{\abovedisplayskip}{4pt plus 1pt minus 1pt}
\setlength{\belowdisplayskip}{4pt plus 1pt minus 1pt}
\setlength{\abovedisplayshortskip}{2pt plus 1pt}
\setlength{\belowdisplayshortskip}{2pt plus 1pt}

\setlength{\textfloatsep}{7pt plus 2pt minus 2pt}
\setlength{\floatsep}{6pt plus 2pt minus 2pt}
\setlength{\intextsep}{6pt plus 2pt minus 2pt}

\renewcommand{\topfraction}{0.95}
\renewcommand{\bottomfraction}{0.90}
\renewcommand{\textfraction}{0.03}
\renewcommand{\floatpagefraction}{0.85}

\setcounter{topnumber}{5}
\setcounter{bottomnumber}{4}
\setcounter{totalnumber}{8}

\setlist{
topsep=2pt,
itemsep=1pt,
parsep=0pt,
partopsep=0pt
}

\titlespacing*{\section}
{0pt}{8pt plus 2pt minus 1pt}{3pt plus 1pt minus 1pt}

\titlespacing*{\subsection}
{0pt}{6pt plus 1pt minus 1pt}{2pt}

\titlespacing*{\subsubsection}
{0pt}{4pt plus 1pt minus 1pt}{1pt}
%==============================================================================
% Title and authors
%==============================================================================
\title{
Graph-RMS: Training-Free and Class-Count-Free Spectral-Spatial Region
Discovery in Hyperspectral Remote-Sensing Images
}

\author[1]{Simhaa T.\,T.}
\author[2]{Steven Yves Le Moan}
\author[3,*]{Shihabudheen K.\,V.}

\affil[1]{
Department of Electrical and Electronics Engineering,
Sri Sivasubramaniya Nadar College of Engineering,
Kalavakkam, Chennai 603110, Tamil Nadu, India
}

\affil[2]{
Department of Computer Science,
Norwegian University of Science and Technology,
Gjøvik, Norway
}

\affil[3]{
Department of Electrical Engineering,
National Institute of Technology Calicut,
NIT Campus P.O., Kozhikode 673601, Kerala, India
}

\affil[*]{
Corresponding author: shihabudheen@nitc.ac.in
}

\date{}

%==============================================================================
% Document
%==============================================================================
\begin{document}

\maketitle

%==============================================================================
% Abstract
%==============================================================================

\begin{abstract}
Hyperspectral images contain detailed spectral and spatial information, but
their analysis is difficult when labelled samples and a reliable semantic
class count are unavailable. Graph-RMS (Graph-based Robust Mean Shift) is a
deterministic, training-free method for discovering coherent
spectral--spatial regions without supplying reference labels or a prescribed
number of groups during scene processing. It constructs a reciprocal local
graph, propagates the spectral representation through damped diffusion, forms
fine modes, and consolidates compatible region prototypes through guarded
endpoint selection. The primary evaluation used eight development scenes and
Trento as an independent holdout for the frozen procedure. Graph-RMS achieved
the highest normalised mutual information on all nine scenes in the
five-method classical comparison. On Trento, it obtained normalised mutual
information of 0.8495 and adjusted Rand index of 0.9447; the corresponding
one-to-one Hungarian overall and balanced accuracies were 0.8533 and 0.6774.
Fixed-default HDBSCAN was lower in all 36 matched scene--metric comparisons,
while Leiden community detection on the same graph remained competitive. A
separate automatic-selection experiment examined whether the principal
scene-dependent settings could be chosen without target-scene labels or the
semantic class count. Across the eight development scenes, automatic
selection increased mean balanced accuracy, normalised mutual information,
and adjusted Rand index, while mean overall accuracy decreased. The frozen
automatic rule also returned a valid operating point in a retrospective
Trento transfer test, although its agreement was lower than that of the
primary Trento configuration. These results position Graph-RMS as a method
for structural region discovery rather than semantic class-count estimation.
\end{abstract}

\noindent\textbf{Keywords:} hyperspectral remote sensing; unsupervised clustering; spectral--spatial analysis; graph-based region discovery; unknown class count

%==============================================================================
% 1. Introduction
%==============================================================================
%==============================================================================
% 1. Introduction
%==============================================================================
\section{Introduction}
\label{sec:introduction}

Hyperspectral images record a reflectance spectrum at each pixel, allowing
materials with similar visible appearance to be distinguished through their
wavelength-dependent responses. This information has supported applications
ranging from vegetation and crop assessment to mineral mapping, environmental
monitoring, and urban land-cover analysis. The same spectral detail also makes
the images difficult to examine. Hundreds of correlated bands must be interpreted
together, scene dimensions can be large, and the spectral response of a material
may vary with illumination, atmospheric effects, acquisition conditions, and
within-class heterogeneity \citep{bioucasdias2013}. Spatial context helps resolve
some of this ambiguity because nearby pixels often belong to the same physical
region, although boundaries, mixed pixels, and small objects prevent proximity
alone from defining a reliable partition \citep{he2018}.

Supervised classification remains the usual route to thematic maps when a
representative set of labelled pixels is available. Such labels may not exist
during the initial examination of a newly acquired scene. Producing them
requires field observations or expert interpretation, and the resulting
reference map may cover only a small part of the image. The analyst may also
lack a defensible estimate of how many distinct regions are present. This is
particularly relevant during exploratory analysis, where the immediate aim is
not yet to assign every pixel to a fixed land-cover legend, but to identify
spectrally and spatially coherent areas that merit inspection. An unsupervised
partition can therefore serve as an intermediate representation of the scene
before semantic interpretation or targeted annotation begins.

Conventional clustering provides a direct starting point for this task, but its
assumptions do not always match the structure of hyperspectral imagery.
Centroid-based methods usually require the number of clusters to be specified
before fitting and can change with initialisation. They also treat pixels as
independent observations unless spatial information is added explicitly.
A partition obtained from spectral vectors alone may consequently contain
isolated pixels and fragmented regions even when its cluster centres are
well separated. Projected clustering and spectral--spatial sparse subspace
methods reduce some of these difficulties by working in lower-dimensional
representations or by introducing local spatial constraints
\citep{mehta2016,zhang2016}. Their output nevertheless depends on choices such
as representation dimension, neighbourhood construction, regularisation, and
the prescribed number of groups.

Graph representations offer a different way to organise a hyperspectral scene.
Pixels form vertices, while weighted edges retain selected spectral or
spectral--spatial relationships. Clustering then depends on the structure of
local connections rather than on distance from a single global centroid
\citep{vonluxburg2007}. This formulation is useful when a class occupies a
non-convex region in feature space or when local continuity matters more than
global compactness. The graph itself remains consequential: dense neighbourhoods
can connect unrelated areas, whereas an overly restrictive graph can split a
continuous region. Mutual-neighbour rules provide one means of retaining
reciprocal local relationships, but low-density pixels may then become weakly
connected or isolated. A practical graph-based method must therefore balance
local selectivity with sufficient connectivity.

Diffusion methods extend this local view by propagating information through
successive graph transitions. The resulting representation reflects paths
through the neighbourhood graph and can reveal relationships that are not
captured by direct pairwise distance \citep{coifman2006}. In hyperspectral
analysis, nonlinear diffusion has been used to combine spectral geometry with
spatial information for clustering and active learning \citep{murphy2019}.
More recent studies have coupled diffusion with volume-based criteria or
superpixel regularisation to obtain coherent scene partitions
\citep{polk2023,cui2024}. Graph contrastive methods instead learn an embedding
before clustering, allowing spectral and spatial information to be integrated
through a trainable representation \citep{guan2024}. Recent hyperspectral
clustering work has extended this direction through lightweight
self-supervised graph clustering, hypergraph contrastive learning,
superpixel-level affinity refinement, flexible multi-anchor graphs, and
dynamic sparse graph autoencoders
\citep{ding2025slcgc,yao2025hypergraph,qin2025evolving,
li2025multianchor,liang2026dsgc}. These approaches reflect the current
emphasis on learned graph representations, spatial aggregation, and scalable
graph construction. Graph-RMS addresses a different operating setting: it
does not train a representation network and does not require a prescribed
semantic class count during target-scene processing.

Methods that do not require an explicit cluster count provide another relevant
comparison. HDBSCAN extracts stable regions from a hierarchy of density-connected
samples and can leave weakly supported observations as noise
\citep{campello2013}. Leiden community detection partitions a graph by
optimising community structure without receiving a target number of communities
\citep{traag2019}. Both avoid direct specification of \(K\), but their outputs
still depend on density, neighbourhood, resolution, and graph-construction
choices. More importantly, a data-derived region count need not coincide with
the number of categories in a reference map. A semantic class may contain
several spectral or spatial regimes, while different semantic classes may
overlap in the measured feature space. Cluster validity must therefore be
considered in relation to the purpose of the partition rather than reduced to
agreement with a single presumed ``true'' grouping \citep{hennig2015}.

Evaluation also depends on the information available to each method. A method
applied only to labelled pixels avoids unlabelled areas that may shape
full-scene connectivity, while a method supplied with the reference value of
\(K\) solves a different problem from one that must determine its partition
without \(K\). We therefore process complete scenes, apply the reference masks
only during scoring, and report the true-\(K\) and class-count-free comparisons
separately.

The present study examines how much spectral--spatial organisation can be
recovered without labelled training samples, learned model parameters, or a
prescribed semantic class count. We introduce Graph-RMS (Graph-based Robust
Mean Shift), a deterministic procedure for full-scene region discovery. The
method is inspired by the repeated local weighted-mean updates underlying
mean-shift mode seeking \citep{cheng1995,comaniciu2002}, but its update is a
finite damped diffusion on a fixed spectral--spatial graph rather than
conventional kernel mean shift. Pixels form graph nodes, reciprocal local
edges encode spectral and spatial similarity, and each diffusion step moves a
representation part-way towards the graph-weighted mean of its neighbours.
Nearby diffused representations form fine modes,
which are spatially refined and consolidated through reciprocal prototype
comparisons. Minimum-degree repair and guarded endpoint selection prevent
weak graph support and degenerate consolidation from being accepted
uncritically. The resulting partition describes coherent spectral--spatial
structure in the complete image rather than attempting to reconstruct a
predetermined semantic legend.

The evaluation uses nine hyperspectral scenes spanning different spatial
scales, spectral resolutions, acquisition platforms, and levels of
reference-label coverage. Eight scenes were available during method
development, and their reference-based evaluations were also available while
the procedure and operating settings were being established. The term
``label-free'' therefore refers to the processing and endpoint selection
performed for an individual scene, rather than to the broader development
history of the method. Reference labels and the semantic class count are not
inputs to graph construction, diffusion, fine-mode formation, prototype
consolidation, or endpoint selection. Trento was excluded from development of
the primary procedure and was processed only after that procedure had been
fixed, providing its independent holdout evaluation. The comparison further
separates four classical methods supplied with the reference value of \(K\)
from HDBSCAN and Leiden, which operate without this information. HDBSCAN is
included as a fixed-default density-based reference, while Leiden is applied
to the same weighted graph as Graph-RMS.

The primary experiments retain scene-specific diffusion, fine-mode, and
prototype-consolidation settings. A separate automatic-selection experiment
was therefore conducted to examine whether these operating choices could be
made without target-scene labels or knowledge of the semantic class count. A
single global rule was calibrated on the eight development scenes and then
fixed. Once calibrated, the rule selected an operating point for an individual
scene using only the hyperspectral cube and partition statistics. The same
rule was subsequently applied to Trento as a retrospective transfer test. This
test is reported separately from the primary Trento holdout because the
original Trento result was already known when the automatic-selection
experiment was designed.

The contribution of this study is a graph-based Robust Mean Shift procedure
that combines reciprocal spectral--spatial graph construction,
minimum-degree repair, damped diffusion, fine-mode formation, and guarded
prototype consolidation without requiring the semantic class count during
scene processing. The study also separates structural
region discovery from semantic class-count estimation through full-scene
evaluation, explicit treatment of unmatched regions, and comparisons with both
true-\(K\) and no-\(K\) methods. Component ablation, parameter sensitivity,
fragmentation controls, perturbation tests, cross-device agreement, spatial
bootstrap analysis, computational profiling, and the additional automatic
operating-point experiment are used to examine where the method succeeds and
where its behaviour changes. The intended output is a set of coherent
spectral--spatial regions for exploratory mapping and pre-annotation analysis,
not an estimate of the number of semantic classes.

%==============================================================================
% 3. Materials and methods
%==============================================================================
%==============================================================================
% 2. Materials and methods
%==============================================================================
\section{Materials and methods}
\label{sec:materials_methods}

%==============================================================================
% 2.1 Problem formulation and information-access protocol
%==============================================================================
\subsection{Problem formulation and information-access protocol}
\label{subsec:problem_formulation}

Let a hyperspectral image be represented by
\(\mathbf{X}\in\mathbb{R}^{H\times W\times B}\), where \(H\) and \(W\)
denote the spatial dimensions and \(B\) is the number of spectral bands. After
vectorisation, the scene contains \(N=HW\) pixels. The spectral response of
pixel \(i\) is denoted by \(\mathbf{x}_i\in\mathbb{R}^{B}\), and its spatial
location by \(\mathbf{p}_i\in\mathbb{R}^{2}\). Graph-RMS uses the paired
spectral and spatial observations

\begin{equation}
\mathcal{D}
=
\left\{
\left(\mathbf{x}_i,\mathbf{p}_i\right)
\right\}_{i=1}^{N}
\end{equation}

to produce a partition of the complete scene,

\begin{equation}
\widehat{\mathcal{P}}
=
\left\{
\widehat{\mathcal{R}}_{1},
\widehat{\mathcal{R}}_{2},
\ldots,
\widehat{\mathcal{R}}_{\widehat{M}}
\right\},
\end{equation}

where \(\widehat{\mathcal{R}}_{m}\) is a discovered spectral--spatial region
and \(\widehat{M}\) is the number of regions retained by the partition-based
selection procedure. Each pixel is assigned a region index
\(\widehat{z}_{i}\in\{1,\ldots,\widehat{M}\}\).

The objective is to identify regions whose pixels are locally consistent in
spectral response and spatial organisation. The task differs from supervised
classification because no labelled pixels are supplied during fitting. It also
differs from conventional clustering protocols that prescribe the number of
output groups in advance. Graph-RMS does not receive the semantic class count
\(K\), and \(\widehat{M}\) is not constrained to reproduce it. A single
reference class may contain several spectrally distinct regions, while similar
spectral responses may occur across different semantic classes. The resulting
partition is therefore interpreted as a structural description of the scene
rather than an estimate of its thematic class count.

For a fixed scene run, reference labels and the semantic class count are
excluded from Graph-RMS processing and endpoint selection. The selected
partition is frozen before comparison with the reference map. Formally, if
\(\mathbf{y}_{\Omega}\) denotes the available reference labels within the
labelled-pixel set \(\Omega\), then

\begin{equation}
\widehat{\mathbf{z}}
=
\mathcal{F}
\left(
\mathbf{X},
\mathbf{P};
\boldsymbol{\theta}
\right),
\qquad
\mathbf{y}_{\Omega}\notin
\mathcal{F},
\qquad
K\notin
\mathcal{F},
\end{equation}

where \(\mathcal{F}\) denotes the Graph-RMS computation for a given scene,
\(\mathbf{P}\) contains the pixel coordinates, and
\(\boldsymbol{\theta}\) contains the operating parameters and partition-based
decision rules. This distinction concerns per-scene execution. Reference-based
results from the eight development scenes were available while the procedure
was being developed and were later used to calibrate the global
automatic-selection rule. They were not inputs to the processing or endpoint
selection of an individual target scene.

Graph construction and region discovery are performed on the full image rather
than only on pixels carrying reference labels. This preserves the spatial and
spectral relationships that extend through unlabelled parts of the scene.
Metrics requiring semantic references are subsequently calculated on
\(\Omega\), whereas full-scene quantities, including the total region count,
fragmentation pattern, runtime, and memory demand, are obtained from the complete
partition. This separation prevents the reference mask from influencing the
regions that are later evaluated.

%==============================================================================
% 2.2 Datasets and holdout design
%==============================================================================
\subsection{Datasets and holdout design}
\label{subsec:datasets}

The experiments used nine public hyperspectral scenes that differ in spatial
extent, spectral dimensionality, land-cover composition, and reference-label
coverage. Salinas-A, Indian Pines, KSC, Pavia University, and Botswana are
established airborne or spaceborne benchmark scenes distributed through the
public hyperspectral remote-sensing scene repository
\citep{gic_hsi_scenes}. The three WHU-Hi scenes---LongKou, HongHu, and
HanChuan---provide UAV-borne imagery of agricultural and rural--urban areas at
high spatial resolution \citep{zhong2020}. Trento contains a rural scene
acquired over southern Trento, Italy \citep{ghamisi2017}. Although the original
Trento release also includes LiDAR data, the present study used only the
hyperspectral cube and its land-cover reference.

Graph-RMS processed every pixel in each scene. The reference masks were not
used to crop the images or determine the graph domain. They were applied only
after the final full-scene partition had been selected. This distinction is
important because the proportion of labelled pixels varies markedly across the
datasets. Labels cover less than 2\% of KSC and Botswana, but more than 86\% of
HongHu and 92\% of LongKou. The full-scene graph can therefore include extensive
areas that do not contribute directly to reference-based accuracy measures.

The first eight scenes in Table~\ref{tab:datasets} formed the development set.
Their reference labels and reference-based metrics were available while the
primary procedure and its operating settings were being established. Trento
was excluded from this development process. After the primary procedure had
been fixed, the Trento cube was processed and the selected partition was
frozen before reference-based evaluation. Trento therefore provides the
independent holdout for the primary experiment. The later automatic-selection
experiment was designed after this holdout result was known, so its Trento
run is treated separately as a retrospective transfer test. The primary
Trento experiment therefore evaluates a frozen, pre-specified scene
configuration; it should not be interpreted as a test of automatic scale
selection.

\begin{table}[htbp]
\centering
\small
\caption{Characteristics and experimental roles of the nine hyperspectral
scenes. Reference classes and labelled pixels describe the available evaluation
maps; Graph-RMS processed each complete scene.}
\label{tab:datasets}

\renewcommand{\arraystretch}{1.15}
\setlength{\tabcolsep}{4.2pt}

\begin{tabular}{@{}lccrrrl@{}}
\toprule
Dataset &
\makecell{Spatial\\size} &
Bands &
\makecell{Reference\\classes} &
\makecell{Labelled\\pixels} &
\makecell{Coverage\\(\%)} &
Role \\
\midrule

Salinas-A
& $83 \times 86$
& 204
& 6
& 5\,348
& 74.92
& Development \\

Indian Pines
& $145 \times 145$
& 200
& 16
& 10\,249
& 48.75
& Development \\

KSC
& $512 \times 614$
& 176
& 13
& 5\,211
& 1.66
& Development \\

Pavia University
& $610 \times 340$
& 103
& 9
& 42\,776
& 20.62
& Development \\

WHU-Hi LongKou
& $550 \times 400$
& 270
& 9
& 204\,542
& 92.97
& Development \\

WHU-Hi HongHu
& $940 \times 475$
& 270
& 22
& 386\,693
& 86.61
& Development \\

WHU-Hi HanChuan
& $1217 \times 303$
& 274
& 16
& 257\,530
& 69.84
& Development \\

Botswana
& $1476 \times 256$
& 145
& 14
& 3\,248
& 0.86
& Development \\

Trento
& $166 \times 600$
& 63
& 6
& 30\,214
& 30.34
& Primary holdout \\

\bottomrule
\end{tabular}
\end{table}


\FloatBarrier

%==============================================================================
% 2.3 Spectral preprocessing and local graph construction
%==============================================================================
%==============================================================================
% 2.3 Spectral preprocessing and local graph construction
%==============================================================================
\subsection{Spectral preprocessing and local graph construction}
\label{subsec:graph_construction}

Graph-RMS represented each complete hyperspectral scene as a sparse weighted
graph in which pixels formed the vertices and local spectral--spatial
relationships formed the edges. For each reported scene run, reference labels
were not inputs to graph construction. The same frozen graph was subsequently
used for diffusion, fine-region extraction, and the matched Leiden comparison.
Figure~\ref{fig:graphrms_workflow} summarises the complete processing sequence.

\begin{figure}[htbp]
\centering
\includegraphics[width=\textwidth]{figures/graph.pdf}
\caption{Processing sequence of Graph-RMS. PCA-20 was used only to rank
candidate neighbours within a local spatial window, whereas edge weighting and
graph diffusion used the complete standardised spectra. Reciprocal neighbour
selection and minimum-degree repair produced the sparse graph support. Damped
diffusion, fine-mode extraction, spatial cleanup, and reciprocal prototype
consolidation generated candidate full-scene partitions, from which a
partition-based rule selected the reported endpoint. Reference labels and the
semantic class count were introduced only after the selected partition had
been frozen.}
\label{fig:graphrms_workflow}
\end{figure}

For pixel \(i\), let
\(\mathbf{x}_i=[x_{i1},\ldots,x_{iB}]^{\mathsf T}\in\mathbb{R}^{B}\)
denote its spectrum and
\(\mathbf{p}_i=[r_i,c_i]^{\mathsf T}\in\mathbb{R}^{2}\) its row and column
coordinates. Each spectral band was standardised across the full scene:

\begin{equation}
\widetilde{x}_{ib}
=
\frac{x_{ib}-\mu_b}{s_b},
\label{eq:band_standardisation}
\end{equation}

where

\begin{equation}
\mu_b
=
\frac{1}{N}
\sum_{i=1}^{N}x_{ib},
\qquad
s_b^{\mathrm{raw}}
=
\sqrt{
\frac{1}{N}
\sum_{i=1}^{N}
\left(x_{ib}-\mu_b\right)^2
},
\label{eq:band_statistics}
\end{equation}

with the numerical guard

\begin{equation}
s_b
=
\begin{cases}
s_b^{\mathrm{raw}}, & s_b^{\mathrm{raw}}\geq10^{-8},\\
1, & s_b^{\mathrm{raw}}<10^{-8}.
\end{cases}
\label{eq:band_std_guard}
\end{equation}

Thus, spectral standardisation used the population standard deviation across
all pixels in the scene. The guard prevented division by a numerically
degenerate band scale. This operation gave each band a comparable numerical
scale and prevented high-variance wavelengths from dominating Euclidean
distance. The complete
standardised spectrum
\(\widetilde{\mathbf{x}}_i\in\mathbb{R}^{B}\) was retained for edge weighting
and graph diffusion.

Principal component analysis (PCA) was used only to reduce the cost of local
neighbour ranking \citep{jolliffe2016}. The standardised scene was projected
onto its first 20 principal components:

\begin{equation}
\mathbf{u}_i
=
\mathbf{V}_{20}^{\mathsf T}
\widetilde{\mathbf{x}}_i,
\qquad
\mathbf{u}_i\in\mathbb{R}^{20},
\label{eq:pca_search}
\end{equation}

where
\(\mathbf{V}_{20}\in\mathbb{R}^{B\times20}\) contains the loading vectors
associated with the 20 largest eigenvalues of the scene-level covariance
matrix. The reduced vectors \(\mathbf{u}_i\) were used only to identify
candidate neighbours. They did not replace the full spectra in the affinity
or diffusion calculations.

Candidate neighbours were restricted to a square spatial window centred on
each pixel. With a window radius of seven pixels, the candidate set was

\begin{equation}
\mathcal{W}_i
=
\left\{
j\neq i:
\left\|
\mathbf{p}_i-\mathbf{p}_j
\right\|_{\infty}
\leq 7
\right\},
\label{eq:local_window}
\end{equation}

which corresponds to a maximum \(15\times15\) neighbourhood for an interior
pixel. At image boundaries, the window was clipped to the available scene
extent. Within \(\mathcal{W}_i\), the Euclidean distances

\begin{equation}
d^{\mathrm{PCA}}_{ij}
=
\left\|
\mathbf{u}_i-\mathbf{u}_j
\right\|_2
\label{eq:pca_distance}
\end{equation}

were ranked, and up to 20 nearest local candidates were retained. The resulting
directed candidate list for pixel \(i\) is denoted by \(\mathcal{N}_i\).

An undirected edge was initially admitted only when the neighbourhood
relationship was reciprocal:

\begin{equation}
M_{ij}
=
\begin{cases}
1,
&
j\in\mathcal{N}_i
\ \text{and}\
i\in\mathcal{N}_j,
\\[3pt]
0,
&
\text{otherwise}.
\end{cases}
\label{eq:mutual_neighbour_gate}
\end{equation}

The reciprocal gate removed one-sided nominations that were not supported by
the local neighbourhood geometry. Its selectivity could, however, leave
pixels in sparse or transitional areas with insufficient connectivity.
Graph-RMS therefore applied a minimum-degree repair. For every vertex with
degree below four, the closest remaining candidates from
\(\mathcal{W}_i\) were added in ascending order of
\(d^{\mathrm{PCA}}_{ij}\) until degree four was reached or no candidate
remained. Each added link was symmetrised, and self-loops were excluded.

Let \(S_{ij}\in\{0,1\}\) denote the final graph support after reciprocal
selection and minimum-degree repair. A supported edge received a product of
spectral and spatial Gaussian affinities:

\begin{equation}
A_{ij}
=
S_{ij}
\exp
\left[
-\frac{
\left\|
\widetilde{\mathbf{x}}_i-
\widetilde{\mathbf{x}}_j
\right\|_2^2
}{
2\sigma_x^2
}
\right]
\exp
\left[
-\frac{
\left\|
\mathbf{p}_i-\mathbf{p}_j
\right\|_2^2
}{
2\sigma_p^2
}
\right].
\label{eq:spectral_spatial_affinity}
\end{equation}

The Gaussian bandwidths were estimated from a deterministic sample of the
stored supported edges. Let \(\mathcal{E}_{s}\) contain
\(\min(50\,000,|\mathcal{E}|)\) supported edges sampled without replacement
using random seed 0. The spectral bandwidth was

\begin{equation}
\sigma_x
=
\max
\left\{
\operatorname{median}_{(i,j)\in\mathcal{E}_{s}}
\left\|
\widetilde{\mathbf{x}}_i-
\widetilde{\mathbf{x}}_j
\right\|_2,
10^{-6}
\right\},
\label{eq:spectral_bandwidth}
\end{equation}

and the spatial bandwidth was

\begin{equation}
\sigma_p
=
\max
\left\{
\operatorname{median}_{(i,j)\in\mathcal{E}_{s}}
\left\|
\mathbf{p}_i-\mathbf{p}_j
\right\|_2,
10^{-6}
\right\}.
\label{eq:spatial_bandwidth}
\end{equation}

Both bandwidths were calculated separately for each scene without labels or
scene-specific accuracy optimisation. The fixed sampling seed made the
bandwidth estimation reproducible while limiting the cost of the median
calculation on large graphs. The multiplicative affinity assigned
large weights only when two connected pixels were similar in both spectral
response and spatial location. The resulting matrix
\(\mathbf{A}\in\mathbb{R}^{N\times N}\) was symmetric, non-negative, and
sparse. It was passed unchanged to the damped diffusion stage described in
Section~\ref{subsec:graph_diffusion}.
%==============================================================================
% 2.4 Damped diffusion on the spectral--spatial graph
%==============================================================================
\subsection{Damped diffusion on the spectral--spatial graph}
\label{subsec:graph_diffusion}

The weighted graph provided a local description of the scene, but a direct
edge represented only one neighbourhood step. Graph-RMS used diffusion to
propagate the complete spectral representation through paths in this graph.
This allowed each pixel to incorporate information from progressively wider
spectral--spatial neighbourhoods while retaining the local geometry encoded by
the sparse affinity matrix \citep{coifman2006,murphy2019}.

Let
\(\mathbf{A}\in\mathbb{R}^{N\times N}\) denote the symmetric affinity matrix
defined in Eq.~\eqref{eq:spectral_spatial_affinity}. The corresponding degree
matrix was

\begin{equation}
\mathbf{D}
=
\operatorname{diag}
\left(
\mathbf{A}\mathbf{1}
\right),
\qquad
D_{ii}
=
\sum_{j=1}^{N}A_{ij},
\label{eq:degree_matrix}
\end{equation}

and the row-normalised transition matrix was

\begin{equation}
\mathbf{W}
=
\mathbf{D}^{-1}\mathbf{A}.
\label{eq:transition_matrix}
\end{equation}

Because the minimum-degree repair removed isolated vertices, every diagonal
entry of \(\mathbf{D}\) was positive. Each row of \(\mathbf{W}\) therefore
summed to one, and \(W_{ij}\) represented the proportion of information
transferred from neighbouring pixel \(j\) to pixel \(i\).

The initial node representation was the complete standardised spectral matrix,

\begin{equation}
\mathbf{Y}^{(0)}
=
\widetilde{\mathbf{X}}
\in
\mathbb{R}^{N\times B}.
\label{eq:diffusion_initialisation}
\end{equation}

Graph-RMS updated this representation through

\begin{equation}
\mathbf{Y}^{(t+1)}
=
(1-\alpha)\mathbf{Y}^{(t)}
+
\alpha\mathbf{W}\mathbf{Y}^{(t)},
\label{eq:damped_diffusion}
\end{equation}

with the damping parameter fixed at

\begin{equation}
\alpha=0.5.
\label{eq:diffusion_alpha}
\end{equation}

Equivalently, the representation of pixel \(i\) at iteration \(t+1\) was

\begin{equation}
\mathbf{y}^{(t+1)}_i
=
\frac{1}{2}\mathbf{y}^{(t)}_i
+
\frac{1}{2}
\sum_{j=1}^{N}
W_{ij}\mathbf{y}^{(t)}_j.
\label{eq:node_diffusion}
\end{equation}

Each update therefore retained half of the current pixel representation and
replaced the remaining half with an affinity-weighted average of its graph
neighbours. The self-retention term damped the transition process and avoided
the two-cycle oscillation that can occur on reciprocal local graph structures.
No model parameters were learned during diffusion. The update is
mean-shift-inspired in its repeated movement towards a local weighted mean, but
it is not mathematically equivalent to conventional mean shift
\citep{cheng1995,comaniciu2002}. Conventional mean shift derives its update
from kernel-weighted neighbourhoods in feature space, whereas Graph-RMS fixes
the graph support and transition weights before diffusion and retains the
representation at a finite checkpoint.

Graph-RMS did not run the update until convergence. Instead, it retained the
representation at a finite checkpoint \(T\),

\begin{equation}
\mathbf{Y}^{(T)}
=
\left[
(1-\alpha)\mathbf{I}
+
\alpha\mathbf{W}
\right]^{T}
\widetilde{\mathbf{X}}.
\label{eq:checkpoint_representation}
\end{equation}

The checkpoint controlled the graph scale represented in
\(\mathbf{Y}^{(T)}\). Small values preserved more of the initial pixel-level
variation, whereas additional iterations incorporated information transmitted
through longer graph paths. The scene-specific checkpoints used in the
primary experiment were established during method development and then frozen
for the reported evaluation.

The reported configuration used \(T=100\) for Salinas-A, KSC, Pavia
University, WHU-Hi LongKou, WHU-Hi HanChuan, and Trento; \(T=50\) for Indian
Pines and Botswana; and \(T=25\) for WHU-Hi HongHu. These values were fixed for the primary experiment before its reported
evaluation. Reference-based results from the development scenes were available
during method development, whereas the Trento reference map was not used to
revise its held-out configuration. The effect of \(T\) was examined separately
in the sensitivity analysis.

The update was implemented using sparse compressed-row matrix multiplication.
Only the non-zero entries of \(\mathbf{A}\) and \(\mathbf{W}\) were stored, so
the calculation operated on the retained local graph rather than forming a
dense \(N\times N\) matrix. The checkpoint representation
\(\mathbf{Y}^{(T)}\) was then passed to the fine-mode extraction stage.

%==============================================================================
% 2.5 Fine-mode formation and spatial refinement
%==============================================================================
\subsection{Fine-mode formation and spatial refinement}
\label{subsec:fine_modes}

The finite diffusion checkpoint retained gradual variation within broad
spectral--spatial structures. Graph-RMS converted this representation into an
initial fine partition before attempting any region consolidation. The
checkpoint matrix \(\mathbf{Y}^{(T)}\) was first projected onto its leading ten
principal components \citep{jolliffe2016}:

\begin{equation}
\mathbf{h}_i
=
\mathbf{V}_{10}^{\mathsf T}
\mathbf{y}_i^{(T)},
\qquad
\mathbf{h}_i\in\mathbb{R}^{10},
\label{eq:fine_pca}
\end{equation}

where \(\mathbf{V}_{10}\) contains the loading vectors associated with the ten
largest eigenvalues of the checkpoint covariance matrix. Each retained
component was then standardised across the scene,

\begin{equation}
z_{i\ell}
=
\frac{h_{i\ell}-\bar{h}_{\ell}}
{s_{h,\ell}},
\qquad
\ell=1,\ldots,10,
\label{eq:fine_standardisation}
\end{equation}

The component standard deviations were calculated as population standard
deviations. If a component standard deviation was below \(10^{-8}\), it was
replaced by 1 before standardisation to avoid division by a numerically
degenerate scale.

The resulting representation was
\(\mathbf{z}_i=[z_{i1},\ldots,z_{i10}]^{\mathsf T}\). This second PCA operation
served a different purpose from PCA-20 in Section~\ref{subsec:graph_construction}.
PCA-20 reduced the cost of the original graph search, whereas PCA-10 described
the geometry of the diffused checkpoint used to form fine regions.

Graph-RMS next constructed a radius graph in the ten-dimensional checkpoint
space. For computational control, no more than 50 nearest candidates were
examined for each pixel. Let \(\mathcal{Q}^{50}_i\) denote this capped candidate
set. An undirected fine-graph edge was retained when a candidate pair satisfied

\begin{equation}
\left\|
\mathbf{z}_i-\mathbf{z}_j
\right\|_2
\leq r_f,
\label{eq:fine_radius_condition}
\end{equation}

where \(r_f\) is the frozen fine-mode radius for the scene. The fine graph can
therefore be written as

\begin{equation}
\mathcal{G}_{f}
=
\left(
\mathcal{V},
\mathcal{E}_{f}
\right),
\qquad
\mathcal{E}_{f}
=
\left\{
\{i,j\}:
j\in\mathcal{Q}^{50}_i
\ \text{or}\
i\in\mathcal{Q}^{50}_j,
\
\left\|
\mathbf{z}_i-\mathbf{z}_j
\right\|_2
\leq r_f
\right\}.
\label{eq:fine_graph}
\end{equation}

The connected components of \(\mathcal{G}_{f}\) defined the initial fine
modes. If the graph contained \(M_f\) components, the corresponding label of
pixel \(i\) was

\begin{equation}
c_i
\in
\left\{
1,\ldots,M_f
\right\}.
\label{eq:fine_labels}
\end{equation}

This construction allowed a fine mode to follow connected structure in the
diffused representation rather than forcing pixels into a prescribed number
of groups.

The frozen values of \(r_f\) were 0.150 for Salinas-A, 0.300 for Indian Pines,
0.0125 for KSC, 0.150 for Pavia University, 0.200 for WHU-Hi LongKou,
0.200 for WHU-Hi HongHu, 0.075 for WHU-Hi HanChuan, 0.150 for Botswana,
and 0.200 for Trento. These values determined the granularity of the fine partition and formed part
of the frozen primary configuration. Development-scene reference results were
available while the operating settings were being established; no
reference-based adjustment was made during the reported scene runs.

Fine-mode labels were subsequently refined using a spatial majority filter.
For pixel \(i\), the spatial neighbourhood was

\begin{equation}
\mathcal{B}^{(3)}_i
=
\left\{
j:
\left\|
\mathbf{p}_i-\mathbf{p}_j
\right\|_{\infty}
\leq 3
\right\},
\label{eq:cleanup_neighbourhood}
\end{equation}

which corresponds to a \(7\times7\) window for an interior pixel. For each
fine-mode label \(a\), its local count was

\begin{equation}
n_{ia}
=
\sum_{j\in\mathcal{B}^{(3)}_i}
\mathbb{I}
\left(
c_j=a
\right).
\label{eq:local_label_count}
\end{equation}

When one label had a unique maximum count, that label was assigned to the
centre pixel. If two or more labels shared the maximum, the centre label was
retained:

\begin{equation}
\widetilde{c}_i
=
\begin{cases}
\displaystyle
\arg\max_a n_{ia},
&
\text{if the maximum is unique},
\\[6pt]
c_i,
&
\text{otherwise}.
\end{cases}
\label{eq:majority_cleanup}
\end{equation}

Retaining the centre label on ties avoided changes caused only by arbitrary
numeric ordering of the region identifiers. The cleaned labels
\(\widetilde{c}_i\) formed the fine-region input to prototype consolidation.

%==============================================================================
% 2.6 Reciprocal consolidation of region prototypes
%==============================================================================
%==============================================================================
% 2.6 Reciprocal consolidation of region prototypes
%==============================================================================
\subsection{Reciprocal consolidation of region prototypes}
\label{subsec:prototype_consolidation}

The cleaned fine partition was deliberately detailed and could retain several
modes within a broader spectral--spatial structure. Graph-RMS consolidated
compatible fine regions by comparing their prototypes in the standardised
diffusion-PCA space. This stage used only the diffused features, region
memberships, and region sizes; reference labels and the semantic class count
remained unavailable.

Let

\begin{equation}
\mathcal{R}_a
=
\left\{
i:\widetilde{c}_i=a
\right\}
\label{eq:fine_region}
\end{equation}

denote fine region \(a\), with size

\begin{equation}
n_a
=
\left|\mathcal{R}_a\right|.
\label{eq:region_size}
\end{equation}

Its prototype centroid was calculated from the standardised diffusion-PCA
features \(\mathbf{z}_i\in\mathbb{R}^{d}\), where \(d\leq10\):

\begin{equation}
\boldsymbol{\mu}_a
=
\frac{1}{n_a}
\sum_{i\in\mathcal{R}_a}
\mathbf{z}_i.
\label{eq:prototype_centroid}
\end{equation}

The raw within-region dispersion was the root-mean-square Euclidean radius,

\begin{equation}
s_a^{\mathrm{raw}}
=
\sqrt{
\frac{1}{n_a}
\sum_{i\in\mathcal{R}_a}
\left\|
\mathbf{z}_i-\boldsymbol{\mu}_a
\right\|_2^2
}.
\label{eq:raw_prototype_dispersion}
\end{equation}
Here, root mean square refers only to the dispersion statistic in
Eq.~\eqref{eq:raw_prototype_dispersion}; the RMS acronym in Graph-RMS denotes
Robust Mean Shift.

This quantity is equivalent to the square root of the trace of the population
covariance matrix of region \(a\), using a divisor of \(n_a\) rather than
\(n_a-1\).

Very small or singleton regions can have zero or near-zero dispersion. To
prevent them from producing unstable normalised distances, Graph-RMS formed
the reference set

\begin{equation}
\mathcal{S}
=
\left\{
s_b^{\mathrm{raw}}:
s_b^{\mathrm{raw}}>10^{-6},
\;
n_b\geq5
\right\}.
\label{eq:dispersion_reference_set}
\end{equation}

The dispersion floor was then defined as

\begin{equation}
f
=
\begin{cases}
Q_{0.25}\!\left(\mathcal{S}\right),
&
\mathcal{S}\neq\varnothing,
\\[4pt]
10^{-3},
&
\mathcal{S}=\varnothing,
\end{cases}
\label{eq:dispersion_floor}
\end{equation}

where \(Q_{0.25}\) denotes the first quartile. The dispersion used in prototype
comparison was

\begin{equation}
s_a
=
\max
\left(
s_a^{\mathrm{raw}},f
\right).
\label{eq:floored_dispersion}
\end{equation}

The size term was defined relative to the median size of the current
prototypes:

\begin{equation}
q_a
=
\frac{
n_a
}{
\max\left(
\operatorname{median}_{b}n_b,1
\right)
}.
\label{eq:prototype_size_ratio}
\end{equation}

Graph-RMS generated a conservative and a size-aware consolidation endpoint
using

\begin{equation}
\beta\in\{0,0.5\}.
\label{eq:size_exponents}
\end{equation}

For a given \(\beta\), the size-adjusted dispersion was

\begin{equation}
\widehat{s}_a^{(\beta)}
=
\frac{s_a}{
\left[
\max\left(q_a,10^{-8}\right)
\right]^{\beta}
}.
\label{eq:size_adjusted_dispersion}
\end{equation}

When \(\beta=0\), prototype comparison depended only on within-region
dispersion. With \(\beta=0.5\), the effective dispersion increased for small
regions, allowing a small region to merge more readily with a compatible
larger neighbour.

The normalised separation between prototypes \(a\) and \(b\) was

\begin{equation}
\delta_{ab}^{(\beta)}
=
\frac{
\left\|
\boldsymbol{\mu}_a-\boldsymbol{\mu}_b
\right\|_2
}{
\max
\left[
\sqrt{
\left(
\widehat{s}_a^{(\beta)}
\right)^2
+
\left(
\widehat{s}_b^{(\beta)}
\right)^2
},
10^{-8}
\right]
}.
\label{eq:prototype_distance}
\end{equation}

The \(10^{-8}\) guard prevented division by zero without materially changing
non-degenerate prototype distances.

At each consolidation round, up to 30 nearest prototype candidates by centroid
distance were retained for every current region. Within this candidate set,
each prototype nominated the candidate having the smallest normalised
separation:

\begin{equation}
b_a^{\star}
=
\arg\min_{b\in\mathcal{Q}^{30}_a}
\delta_{ab}^{(\beta)},
\label{eq:preferred_prototype}
\end{equation}

where \(\mathcal{Q}^{30}_a\) denotes the retained prototype-neighbour set of
region \(a\). A pair \((a,b)\) was eligible for merging only when the
nominations were mutual and the corresponding normalised separation satisfied
the frozen threshold:

\begin{equation}
b=b_a^{\star},
\qquad
a=b_b^{\star},
\qquad
\delta_{ab}^{(\beta)}\leq\tau.
\label{eq:reciprocal_merge_rule}
\end{equation}

Thus, membership in the reciprocal candidate sets alone was not sufficient
for merging; the two prototypes had to be mutually preferred under the
normalised separation. This restriction reduced unrestricted chaining through
the prototype graph.

When regions \(a\) and \(b\) were merged into region \(c\), their size and
centroid were updated exactly:

\begin{equation}
n_c
=
n_a+n_b,
\label{eq:merged_region_size}
\end{equation}

\begin{equation}
\boldsymbol{\mu}_c
=
\frac{
n_a\boldsymbol{\mu}_a+
n_b\boldsymbol{\mu}_b
}{
n_a+n_b
}.
\label{eq:merged_centroid}
\end{equation}

The merged dispersion was updated from the current prototype statistics:

\begin{equation}
s_c
=
\sqrt{
\frac{
n_a
\left[
s_a^2+
\left\|
\boldsymbol{\mu}_a-\boldsymbol{\mu}_c
\right\|_2^2
\right]
+
n_b
\left[
s_b^2+
\left\|
\boldsymbol{\mu}_b-\boldsymbol{\mu}_c
\right\|_2^2
\right]
}{
n_c
}
}.
\label{eq:pooled_dispersion}
\end{equation}

The dispersions entering this recursion already included the floor defined in
Eq.~\eqref{eq:dispersion_floor}. The floor was calculated once from the
initial fine modes and was not re-estimated after subsequent merges. The
update accounts for both within-region dispersion and separation between the
two prototype centroids.

After each round, Graph-RMS recomputed the current prototype neighbourhoods,
normalised separations, and eligible reciprocal pairs. This round-wise update
prevented unrestricted single-link chaining based on distances calculated
before earlier merges. Consolidation ended when a complete round produced no
eligible pair.

The prototype thresholds associated with the reported primary endpoints were
\(\tau=0.75\) for Salinas-A, \(0.60\) for Indian Pines, \(3.00\) for KSC,
\(3.00\) for Pavia University, \(0.50\) for WHU-Hi LongKou, \(0.75\) for
WHU-Hi HongHu, \(3.00\) for WHU-Hi HanChuan, \(0.65\) for Botswana, and
\(0.75\) for Trento. These scene-specific thresholds formed part of the
frozen primary configuration. For a fixed scene run, the conservative and
size-aware branches were then evaluated through the same partition-based
endpoint rule. Eight scenes retained the size-aware endpoint with
\(\beta=0.5\), while WHU-Hi LongKou used the conservative fallback with
\(\beta=0\). Reference labels and reference-based metrics were not inputs to
this branch-selection step.


%==============================================================================
% 2.7 Label-free endpoint selection and abstention
%==============================================================================
\subsection{Partition-based endpoint selection and abstention}
\label{subsec:endpoint_selection}

Prototype consolidation produced a sequence of candidate partitions rather
than a single automatically accepted result. Increasing the consolidation
threshold generally reduced the number of regions, but excessive consolidation
could also create a dominant region or an unbalanced partition. Graph-RMS
therefore selected an endpoint from full-scene partition properties alone.
Reference labels, reference-based metrics, and the semantic class count were
not available to this procedure.

Let

\begin{equation}
\mathcal{P}_{\ell}
=
\left\{
\mathcal{R}_{\ell 1},
\ldots,
\mathcal{R}_{\ell M_{\ell}}
\right\}
\label{eq:candidate_partition}
\end{equation}

denote the partition obtained at consolidation threshold
\(\tau_{\ell}\), where \(M_{\ell}\) is its full-scene region count.
Thresholds were evaluated in increasing order on a fixed candidate grid.

The persistence of a candidate was measured through the adjusted Rand
agreement between consecutive full-scene partitions
\citep{hubert1985}:

\begin{equation}
S_{\ell}
=
\operatorname{ARI}
\left(
\mathcal{P}_{\ell-1},
\mathcal{P}_{\ell}
\right).
\label{eq:selector_stability}
\end{equation}

This use of ARI did not involve a reference map. It compared two candidate
Graph-RMS partitions over the same pixels and measured whether a change in
threshold substantially altered the region assignment.

Three additional quantities were calculated for each candidate. The fraction
of pixels assigned to its largest region was

\begin{equation}
L_{\ell}
=
\max_{m\in\{1,\ldots,M_{\ell}\}}
\frac{
\left|
\mathcal{R}_{\ell m}
\right|
}{N}.
\label{eq:largest_region_fraction}
\end{equation}

Let

\begin{equation}
q_{\ell m}
=
\frac{
\left|
\mathcal{R}_{\ell m}
\right|
}{N}
\label{eq:region_proportion}
\end{equation}

denote the relative size of region \(m\). The normalised partition entropy was

\begin{equation}
E_{\ell}
=
-\frac{
\displaystyle
\sum_{m=1}^{M_{\ell}}
q_{\ell m}\log q_{\ell m}
}{
\log M_{\ell}
},
\qquad
M_{\ell}>1,
\label{eq:normalized_partition_entropy}
\end{equation}

with \(E_{\ell}=0\) when \(M_{\ell}=1\). Higher values indicate that the
partition mass is distributed across several regions rather than concentrated
in one dominant component.

Compression relative to the cleaned fine partition was defined as

\begin{equation}
C_{\ell}
=
1-
\frac{M_{\ell}}{M_f},
\label{eq:selector_compression}
\end{equation}

where \(M_f\) is the number of fine regions before prototype consolidation.
A value of \(C_{\ell}=0\) indicates that consolidation retained the fine-region
count, whereas larger values indicate progressively stronger compression.

A size-aware candidate was admissible only when all four conditions were met:

\begin{equation}
S_{\ell}\geq0.95,
\qquad
L_{\ell}\leq0.40,
\qquad
E_{\ell}\geq0.35,
\qquad
C_{\ell}\geq0.20.
\label{eq:selector_conditions}
\end{equation}

The stability condition required the candidate to lie within a persistent part
of the consolidation sequence. The largest-region and entropy conditions
rejected collapsed or strongly unbalanced partitions, while the compression
condition required consolidation to produce a material reduction in the
fine-region count.

Among the persistent size-aware candidates satisfying
Eq.~\eqref{eq:selector_conditions}, Graph-RMS retained the most compressed
endpoint:

\begin{equation}
\ell^{\star}
=
\arg\min_{\ell\in\mathcal{A}}
M_{\ell},
\label{eq:selected_endpoint}
\end{equation}

where \(\mathcal{A}\) is the set of admissible endpoints. If the size-aware
branch did not contain an admissible endpoint, the procedure retained the
corresponding conservative partition obtained with \(\beta=0\). This fallback
prevented the size term from forcing consolidation when the resulting
partition failed the stability or degeneracy checks.

For adaptive re-selection under perturbation, Graph-RMS could also abstain.
Abstention occurred when neither branch contained a persistent,
non-degenerate endpoint. In that case, the threshold-zero partition was
retained for diagnostic inspection, but the method did not present it as a
selected compressed solution. The selector status---size-aware selection,
conservative fallback, or abstention---was stored with the partition.

Once the endpoint had been selected, its region labels, threshold, size
exponent, and selector status were frozen. Only then was the reference mask
introduced for quantitative evaluation.

\begin{algorithm}[!t]
\caption{Graph-RMS full-scene region discovery}
\label{alg:graphrms}
\small

\begin{algorithmic}

\Require
\parbox[t]{0.82\linewidth}{
Hyperspectral cube \(\mathbf{X}\), pixel coordinates \(\mathbf{P}\),
spatial radius \(r_w\), neighbour count \(k\), minimum degree
\(d_{\min}\), diffusion checkpoint \(T\), fine-mode radius \(r_f\),
and consolidation settings
}

\Ensure
\parbox[t]{0.82\linewidth}{
Selected full-scene partition \(\widehat{\mathcal{P}}\) and endpoint status
}

\Statex
\Statex \textbf{Stage 1: Spectral--spatial graph construction}
\State Standardise each spectral band over the complete scene
\State Project the spectra to PCA-\(20\) for local neighbour ranking
\State Find up to \(k\) candidates within the \(r_w\)-pixel spatial window
\State Retain reciprocal neighbours and repair vertices below \(d_{\min}\)
\State Compute full-spectrum spectral--spatial affinities
\State Form the sparse row-normalised transition matrix \(\mathbf{W}\)

\Statex
\Statex \textbf{Stage 2: Diffusion and fine-mode formation}
\State Initialise \(\mathbf{Y}^{(0)}\) with the standardised full spectra
\For{\(t=0,\ldots,T-1\)}
    \State
    \(\mathbf{Y}^{(t+1)}
    \gets
    (1-\alpha)\mathbf{Y}^{(t)}
    +
    \alpha\mathbf{W}\mathbf{Y}^{(t)}\)
\EndFor
\State Project \(\mathbf{Y}^{(T)}\) to at most ten standardised components
\State Form a capped radius graph using \(r_f\)
\State Assign connected components as fine modes
\State Apply spatial majority cleanup

\Statex
\Statex \textbf{Stage 3: Reciprocal prototype consolidation}
\For{\(\beta\in\{0,0.5\}\)}
    \State Compute region centroids, sizes, and floored dispersions
    \Repeat
        \State Find reciprocal prototype pairs satisfying the merge threshold
        \State Merge eligible pairs and update pooled prototype statistics
    \Until{no eligible pair remains}
    \State Store the resulting candidate partition sequence
\EndFor

\Statex
\Statex \textbf{Stage 4: Partition-based endpoint selection}
\State Evaluate persistence, largest-region fraction, entropy, and compression
\State Select an admissible endpoint, use the conservative fallback, or abstain
\State Freeze \(\widehat{\mathcal{P}}\) before reference-based evaluation

\Statex
\Statex \textit{Reference labels and the semantic class count are not inputs to the steps above.}

\end{algorithmic}
\end{algorithm}

\begin{table}[!t]
\centering
\footnotesize
\caption{Frozen operating settings for the primary Graph-RMS experiment.
Panel A lists settings shared by all scenes, and Panel B gives the
scene-specific diffusion, fine-mode, and consolidation settings.}
\label{tab:frozen_pipeline}

\renewcommand{\arraystretch}{1.05}
\setlength{\tabcolsep}{4.5pt}

\begin{tabular}{@{}ll@{}}
\toprule
\multicolumn{2}{@{}l}{\textit{Panel A: Shared settings}} \\
\addlinespace[2pt]
Quantity & Setting \\
\midrule
Neighbour-ranking representation & PCA-20 \\
Spatial-window radius & \(r_w=7\) pixels \\
Maximum local neighbours & \(k=20\) \\
Minimum repaired degree & \(d_{\min}=4\) \\
Diffusion damping & \(\alpha=0.5\) \\
Fine-mode representation & At most 10 standardised PCs \\
Fine-neighbour cap & 50 \\
Spatial-cleanup radius & 3 pixels \\
Prototype-neighbour cap & 30 \\
Dispersion reference size & \(n_a\geq5\) \\
Dispersion-floor quantile & \(Q_{0.25}\) \\
Dispersion fallback & \(10^{-3}\) \\
Size exponents & \(\beta\in\{0,0.5\}\) \\
Endpoint stability & \(S_{\ell}\geq0.95\) \\
Largest-region fraction & \(L_{\ell}\leq0.40\) \\
Normalised entropy & \(E_{\ell}\geq0.35\) \\
Minimum compression & \(C_{\ell}\geq0.20\) \\
\midrule
\multicolumn{2}{@{}l}{\textit{Panel B: Scene-specific primary settings}} \\
\addlinespace[2pt]
Dataset & \(T,\ r_f,\ \tau,\ \beta\) \\
\midrule
Salinas-A & \(100,\ 0.150,\ 0.75,\ 0.5\) \\
Indian Pines & \(50,\ 0.300,\ 0.60,\ 0.5\) \\
KSC & \(100,\ 0.0125,\ 3.00,\ 0.5\) \\
Pavia University & \(100,\ 0.150,\ 3.00,\ 0.5\) \\
WHU-Hi LongKou & \(100,\ 0.200,\ 0.50,\ 0.0\) \\
WHU-Hi HongHu & \(25,\ 0.200,\ 0.75,\ 0.5\) \\
WHU-Hi HanChuan & \(100,\ 0.075,\ 3.00,\ 0.5\) \\
Botswana & \(50,\ 0.150,\ 0.65,\ 0.5\) \\
Trento & \(100,\ 0.200,\ 0.75,\ 0.5\) \\
\bottomrule
\end{tabular}
\end{table}

\subsection{Automatic operating-point selection}
\label{subsec:auto_selector}

The primary experiment retained scene-specific values of the diffusion
checkpoint, fine-mode radius, and prototype-consolidation settings. A separate
experiment examined whether these operating choices could instead be selected
for an individual scene without using its reference labels or semantic class
count.

The experiment separated global rule calibration from per-scene selection.
One selection rule was calibrated on the eight development scenes from a
predefined rule grid. Reference-based NMI and ARI from these scenes were used
only at this calibration stage to choose the global rule. The selected rule
was then fixed. Subsequent application to an individual scene used only the
hyperspectral cube and partition statistics; its reference labels and \(K\)
were not available to the selector.

Four diffusion checkpoints were examined:

\begin{equation}
T\in\{25,50,100,200\}.
\label{eq:auto_T_grid}
\end{equation}

At checkpoint \(T\), the tenth non-self-neighbour distance was calculated for
each pixel in the standardised diffusion-PCA representation. The scene scale
was defined as

\begin{equation}
s_T
=
\operatorname{median}_{i}
d_{i,(10)},
\label{eq:auto_scene_scale}
\end{equation}

using positive finite distances only. The fine-mode radius was then
normalised to this scale:

\begin{equation}
r_f
=
\gamma s_T,
\qquad
\gamma
\in
\{0.75,1.00,1.25,1.50\}.
\label{eq:auto_radius}
\end{equation}

Prototype consolidation was evaluated over

\begin{equation}
\tau\in
\{
0,0.05,0.10,0.15,0.20,0.25,0.30,0.35,0.40,0.45,
0.50,0.55,0.60,0.65,0.70,0.75,1.0,1.5,2.0,3.0
\},
\label{eq:auto_tau_grid}
\end{equation}

with

\begin{equation}
\beta\in\{0,0.5\}.
\label{eq:auto_beta_grid}
\end{equation}

Candidate endpoints first had to satisfy the existing largest-region, entropy,
compression, and adjacent-threshold persistence requirements. The global rule
search then considered adjacent-\(T/\gamma\) stability thresholds from 0.80
to 0.90, maximum compression values of 0.85, 0.90, and 0.95, and minimum
agreement between the two \(\beta\) branches from 0.50 to 0.90. Rule
selection first maximised the number of development scenes for which an
admissible operating point was found and then maximised mean
\((\mathrm{NMI}+\mathrm{ARI})/2\). Remaining ties favoured higher stability,
lower maximum compression, and stronger agreement between the two
\(\beta\) branches.

The locked rule required adjacent-\(T/\gamma\) stability of at least 0.80,
maximum compression of 0.85, and agreement between the two \(\beta\) branches
of at least 0.90 before the size-aware branch could be selected. Stable support
had to extend across at least two adjacent \(T/\gamma\) settings. Within each
branch, the stable component with the greatest support was retained, with ties
resolved by mean stability. Its representative then favoured higher stability,
proximity to the component centre, and lower compression. The size-aware
branch was retained only when it had greater stable support, or equal support
with at least 0.01 higher mean stability; otherwise the conservative
\(\beta=0\) branch was used. If no stable component remained, the selector
abstained rather than returning an unsupported operating point.

The rule was fixed after calibration on the eight development scenes. Because
the primary Trento result was already known when this experiment was designed,
the later automatic Trento run was treated as a retrospective transfer test
rather than an independent holdout.
%==============================================================================
% 2.8 Comparison methods and fairness controls
%==============================================================================
\subsection{Comparison methods and fairness controls}
\label{subsec:baselines}

The comparison methods were separated according to whether they received the
reference class count. This distinction is necessary because a method supplied
with the true number of classes does not operate under the same information
constraints as Graph-RMS. All methods processed the complete hyperspectral
scene. Reference labels were introduced only after a partition had been
obtained.

The conventional comparison included PCA-KMeans, MiniBatch-KMeans, fuzzy
\(c\)-means (FCM), and SLIC-KMeans. K-means minimises within-cluster
squared distance to a set of centroids \citep{lloyd1982}, while
MiniBatch-KMeans approximates this optimisation through successive
mini-batches \citep{sculley2010}. FCM assigns graded memberships during
optimisation \citep{bezdek1984}. SLIC-KMeans first forms spatially compact
superpixels using SLIC \citep{achanta2012} and then groups their spectral
representations using K-means.

For these four methods, each band was standardised across the full scene and
the spectra were projected onto 20 principal components. The reference class
count \(K\) was supplied as the requested number of clusters. No labelled
pixels were used to estimate the PCA representation or tune the remaining
settings. Each method was evaluated with five fixed random seeds, and the
scene-level mean was used in the classical comparison.

HDBSCAN and Leiden formed a separate comparison because neither method received
the reference class count. HDBSCAN was applied to the standardised PCA-20 representation using a
fixed-default configuration \citep{campello2013}. It was included as a
density-based reference rather than as an exhaustively optimised competitor;
no scene-specific or reference-based parameter tuning was performed. Pixels
identified as noise retained the label \(-1\) as a separate predicted label
and were passed unchanged to the same evaluation procedure. Noise pixels were
therefore not discarded. For OA and BA, the \(-1\) label participated as one
predicted region in the one-to-one Hungarian assignment; for NMI and ARI, it
remained a distinct predicted label.

Leiden community detection was applied to the exact weighted graph constructed
for Graph-RMS \citep{traag2019}. Its resolution was selected from a fixed grid
using only partition-level information, including agreement across selector
runs, agreement between adjacent resolutions, the largest-community fraction,
and normalised partition entropy. Reference labels and the reference class
count were excluded from this selection. The selected resolution was then
evaluated over five fixed seeds.

The Leiden comparison holds graph construction constant. Differences between
Leiden and Graph-RMS therefore arise after the weighted graph has been formed:
Leiden directly optimises graph communities, whereas Graph-RMS applies damped
diffusion, fine-mode formation, spatial refinement, and reciprocal prototype
consolidation. HDBSCAN instead provides a generic density-based comparison on
the PCA representation.

Results from the four true-\(K\) methods and the two no-\(K\) methods were
reported in separate comparison groups. The classical methods were analysed
within the predefined Friedman and Holm-corrected family. HDBSCAN and Leiden
were examined in a separate paired no-\(K\) analysis so that methods operating
under different information conditions were not combined in the same
inferential family.

Recent self-supervised hyperspectral clustering methods learn spectral--spatial
representations from augmented image patches before applying a separate
clustering stage. They were not included in the present comparison because
their training procedures, patch-sampling protocols, augmentation choices, and
frequent use of a prescribed output count differ materially from the
training-free, full-scene, unknown-\(K\) setting examined here. A fair
comparison would require a separate study in which these design choices and
computational budgets are standardised rather than adapting one implementation
to the present protocol.
%==============================================================================
% 2.9 Evaluation measures and statistical analysis
%==============================================================================
\subsection{Evaluation measures and statistical analysis}
\label{subsec:evaluation}

Graph-RMS produced a partition of every pixel in the complete scene, whereas
reference-based evaluation was restricted to the labelled-pixel set
\(\Omega\). Let \(N_{\Omega}=|\Omega|\), let \(y_i\in\{1,\ldots,K\}\)
denote the reference class of labelled pixel \(i\), and let
\(\widehat{z}_i\in\{1,\ldots,\widehat{M}\}\) denote its predicted region.
Here, \(K\) is used only during evaluation and does not enter Graph-RMS fitting
or endpoint selection.

Overall accuracy (OA) and balanced accuracy (BA) require semantic
correspondence between predicted regions and reference classes. This
correspondence was obtained through one-to-one Hungarian assignment
\citep{munkres1957}. Let

\begin{equation}
C_{mk}
=
\sum_{i\in\Omega}
\mathbb{I}
\left(
\widehat{z}_i=m
\ \land\
y_i=k
\right)
\label{eq:contingency_matrix}
\end{equation}

denote the number of labelled pixels shared by predicted region \(m\) and
reference class \(k\). Let \(\mathcal{A}\) denote the set of one-to-one
partial mappings from predicted regions to reference classes. The evaluation
mapping was

\begin{equation}
\pi^{\star}
=
\arg\max_{\pi\in\mathcal{A}}
\sum_{m\in\operatorname{dom}(\pi)}
C_{m,\pi(m)}.
\label{eq:hungarian_mapping}
\end{equation}

Thus, \(\pi^{\star}\) maps a matched predicted region to at most one reference
class, and each reference class can receive at most one predicted region.
When the predicted partition contained more regions than reference classes,
unmatched predicted regions remained unmatched and their pixels were counted
as errors. This treatment penalised semantic fragmentation rather than mapping
several predicted regions to the same reference class.

After alignment, OA was calculated as

\begin{equation}
\mathrm{OA}
=
\frac{1}{N_{\Omega}}
\sum_{i\in\Omega}
\mathbb{I}
\left(
\pi^{\star}(\widehat{z}_i)=y_i
\right),
\label{eq:overall_accuracy}
\end{equation}

where the indicator is zero when \(\widehat{z}_i\) belongs to an unmatched
predicted region. BA was the unweighted mean of the class-specific recalls:

\begin{equation}
\mathrm{BA}
=
\frac{1}{K}
\sum_{k=1}^{K}
\frac{
\displaystyle
\sum_{i\in\Omega}
\mathbb{I}
\left(
y_i=k
\ \land\
\pi^{\star}(\widehat{z}_i)=k
\right)
}{
\displaystyle
\sum_{i\in\Omega}
\mathbb{I}(y_i=k)
}.
\label{eq:balanced_accuracy}
\end{equation}

OA describes the proportion of correctly aligned labelled pixels, whereas BA
gives equal weight to every reference class and is less dominated by large
classes.

Normalised mutual information (NMI) measured the information shared by the
predicted and reference partitions, and the adjusted Rand index (ARI) measured
pairwise partition agreement after correction for chance
\citep{hubert1985,vinh2010}. Both measures are invariant to the numeric region
identifiers and therefore do not require Hungarian alignment.

Because Graph-RMS is designed for structural region discovery rather than a
one-region-per-class semantic partition, NMI and ARI were treated as the
principal measures of partition correspondence. OA and BA were retained as
complementary semantic-alignment measures under one-to-one Hungarian matching.
Their interpretation therefore includes the penalty incurred when several
discovered regions correspond to the same reference class.

Adjusted mutual information (AMI), homogeneity, completeness, and V-measure
were calculated as supplementary information-based diagnostics
\citep{vinh2010}. AMI was included to examine whether NMI was affected by the
number of predicted regions. Homogeneity measures whether individual predicted
regions contain pixels from a single reference class, while completeness
measures whether pixels belonging to one reference class remain within the
same predicted region. V-measure is their harmonic mean.

The full-scene region count,

\begin{equation}
\widehat{M}_{\mathrm{full}}
=
\left|
\left\{
\widehat{z}_i:
i=1,\ldots,N
\right\}
\right|,
\label{eq:full_region_count}
\end{equation}

was distinguished from the number of predicted regions intersecting the
reference mask,

\begin{equation}
\widehat{M}_{\Omega}
=
\left|
\left\{
\widehat{z}_i:
i\in\Omega
\right\}
\right|.
\label{eq:evaluation_region_count}
\end{equation}

This distinction was necessary for scenes such as KSC and Botswana, where only
a small fraction of the complete image carries reference labels.

Semantic fragmentation was examined through the ratios
\(\widehat{M}_{\mathrm{full}}/K\) and
\(\widehat{M}_{\Omega}/K\), many-to-one cluster purity, the difference between
many-to-one purity and one-to-one OA, and the number of predicted regions
required to cover 90\% of each reference class. The dominant-region fraction
and normalised within-class fragmentation entropy were also recorded. These
quantities describe how a structurally detailed partition relates to the
broader semantic categories in the reference map.

Graph-RMS is deterministic under a fixed computational path. Its uncertainty
was therefore not represented by an artificial across-seed standard deviation.
Instead, 300 spatial block-bootstrap replicates were generated for every scene
using the frozen scene-specific block sizes. OA, BA, NMI, and ARI were
recalculated for each replicate, and the 2.5th and 97.5th percentiles formed
the reported 95\% intervals. Spatial blocks were used in place of independent
pixel resampling to preserve part of the local dependence present in image
data \citep{lahiri2003}.

Cross-scene statistical tests were used as supporting summaries rather than as
the principal basis for method comparison. Graph-RMS, PCA-KMeans,
MiniBatch-KMeans, FCM, and SLIC-KMeans were ranked separately for OA, BA,
NMI, and ARI across the nine paired scenes. For each classical baseline, the
mean of five seeded runs represented its scene-level performance. All ranks
and statistical tests were calculated from the full-precision scene-level
values stored in the reproducibility archive; values were rounded only for
presentation in the manuscript. Average ranks and Friedman tests were reported
for each metric \citep{friedman1937}. Exact two-sided Wilcoxon signed-rank tests were also
calculated between Graph-RMS and each classical baseline
\citep{wilcoxon1945}, with Holm's procedure controlling the family-wise error
rate across the 16 pairwise comparisons \citep{holm1979}. With only nine
paired scenes, these tests have limited power and were interpreted alongside
the scene-level values and ranks rather than as evidence required to establish
superiority.

The matched no-\(K\) comparisons were summarised separately. Graph-RMS was
compared with HDBSCAN and Leiden using exact two-sided paired Wilcoxon tests,
with one aggregate value per scene and metric. These values are reported as
supporting cross-scene summaries because the no-\(K\) methods address a
different information-access setting and the sample of benchmark scenes is
small.

%==============================================================================
% 2.10 Diagnostic analyses and computational profiling
%==============================================================================
\subsection{Diagnostic analyses and computational assessment}
\label{subsec:additional_analysis}

The primary comparison was supplemented by experiments designed to examine
where the Graph-RMS partition came from, how it changed around the frozen
settings, and what computational resources were required. These analyses were
specified separately from the principal nine-scene comparison and were not
used to modify the Trento holdout configuration.

\subsubsection{Component ablation}

Component ablation was conducted on Pavia University because the scene is large
enough to expose graph and consolidation behaviour while remaining practical
for repeated execution. Each variant altered one part of the frozen procedure
and retained the remaining settings.

The evaluated variants were:

\begin{enumerate}[label=(\roman*)]
\item removal of the mutual-neighbour gate;
\item removal of the spatial term from the graph affinity;
\item removal of diffusion by setting \(T=0\);
\item removal of the spatial majority cleanup;
\item removal of fine-mode grouping;
\item complete removal of prototype consolidation; and
\item consolidation without size weighting, corresponding to
      \(\beta=0\).
\end{enumerate}

For every variant, the full-scene region count, OA, BA, NMI, ARI, and
end-to-end runtime were recorded. The ablation was interpreted component by
component because removal of one stage can improve one semantic metric while
reducing another. No single metric was used as a sufficient test of component
necessity.

\subsubsection{Parameter sensitivity}

Sensitivity was assessed on the development scenes Salinas-A, Indian Pines,
and Pavia University. Four operating parameters were varied one at a time:
the local-window radius, the PCA dimension used for neighbour ranking, the
diffusion checkpoint, and the fine-mode radius. All parameters not under study
were held at their frozen scene values.

The tested window radii were

\begin{equation}
r_w\in\{3,5,7,9\},
\label{eq:sensitivity_window}
\end{equation}

and the tested PCA dimensions were

\begin{equation}
d_{\mathrm{PCA}}
\in
\{10,20,30,40\}.
\label{eq:sensitivity_pca}
\end{equation}

For Salinas-A and Pavia University, the diffusion checkpoints were

\begin{equation}
T\in\{25,50,100,200\},
\label{eq:sensitivity_checkpoint_sp}
\end{equation}

whereas Indian Pines used

\begin{equation}
T\in\{12,25,50,100\}.
\label{eq:sensitivity_checkpoint_ip}
\end{equation}

The fine-mode radii for Salinas-A and Pavia University were

\begin{equation}
r_f
\in
\{0.075,0.1125,0.150,0.1875,0.225\},
\label{eq:sensitivity_radius_sp}
\end{equation}

and Indian Pines used

\begin{equation}
r_f
\in
\{0.150,0.225,0.300,0.375,0.450\}.
\label{eq:sensitivity_radius_ip}
\end{equation}

The complete sensitivity study contained 51 runs. Each run recorded the final
full-scene region count, largest-region fraction, normalised partition entropy,
OA, BA, NMI, ARI, runtime, and memory use. Single-region outputs and other
selector-degenerate outcomes were identified explicitly rather than removed
from the sensitivity range.

\subsubsection{Endpoint-selector audit}

The partition-based selector was examined independently of the primary
performance comparison. For every candidate consolidation endpoint, the stored
reference-free partition statistics were compared with the reference-based
metrics only after selection had been completed. This analysis measured the gap between the
selected endpoint and the best candidate that could have been chosen
retrospectively for OA, BA, NMI, or ARI. These oracle gaps were diagnostic and
did not alter the reported partition.

Rule ablations also examined the consequences of removing the persistence,
largest-region, entropy, and compression requirements. In particular, a
stability-only alternative was tested to determine whether high agreement
between adjacent consolidation thresholds was sufficient to prevent a
collapsed partition.

\subsubsection{Perturbation response and repeatability}

Perturbation experiments were conducted on Salinas-A, Indian Pines, and Pavia
University. Two independently seeded trials were performed for additive noise
at the frozen strength of 0.01 and for random pixel replacement at the frozen
fraction of 0.05. The same preprocessing, graph, diffusion, and consolidation
procedure was rerun on each perturbed scene.

Salinas-A and Indian Pines were initially assessed using their frozen original
consolidation thresholds. For Pavia University, the full partition-based endpoint
rule was rerun for every perturbation because retaining the original threshold
could create a collapse that the adaptive selector would not accept. The
Pavia analysis therefore recorded whether the selector returned a valid
endpoint, retained the conservative candidate, or abstained.

Perturbed partitions were compared with the corresponding frozen partition
using ARI over all scene pixels and, separately, over the labelled support.
Reference-based OA, BA, NMI, and ARI were also calculated after the perturbed
partition had been selected. This separated partition repeatability from
semantic agreement with the reference map.

Cross-device repeatability was examined on KSC, WHU-Hi HanChuan, and Botswana.
The frozen CPU partition was compared with a new GPU execution using full-scene
and labelled-support ARI. Region counts were reported alongside these
agreements because two highly similar partitions can still differ in the
number of small regions. The assessment therefore concerned partition-level
repeatability rather than bitwise identity.

\subsubsection{Computational complexity}

Let \(N\) denote the number of pixels, \(B\) the number of spectral bands,
\(E\) the number of retained graph edges, and \(T\) the diffusion checkpoint.
The local graph examines at most \(w=15^{2}-1\) spatial candidates per pixel
and retains \(k=20\) neighbours. Candidate ranking in the 20-dimensional
search representation therefore requires

\begin{equation}
O\!\left(Nwd_s\right),
\qquad d_s=20,
\end{equation}

while full-spectrum edge weighting requires \(O(EB)\). Since the fixed-degree
construction gives \(E=O(Nk)\), the graph remains sparse.

Each diffusion update multiplies the sparse transition matrix by the
\(N\times B\) spectral representation and therefore costs \(O(EB)\). The full
diffusion stage has complexity

\begin{equation}
O\!\left(TEB\right)
=
O\!\left(TNkB\right).
\label{eq:diffusion_complexity}
\end{equation}

The current implementation identifies up to \(q=50\) fine-mode neighbours
through an exact global nearest-neighbour search in the
\(d_f\leq10\) diffusion-PCA representation. The GPU implementation evaluates
the distance matrix in tiles but compares each query pixel with all \(N\)
pixels. The resulting search therefore has time complexity

\begin{equation}
O\!\left(N^{2}d_f\right).
\label{eq:fine_knn_complexity}
\end{equation}

Only the \(q\) nearest neighbours are retained for each pixel. Once these
neighbours have been identified, radius filtering, connected-component
extraction, and spatial cleanup require \(O(Nq)\) time. Forming the initial
region prototypes costs \(O(Nd_f)\).

Let \(M\) denote the number of current prototypes, \(p=30\) the retained
prototype-neighbour count, and \(R\) the number of consolidation rounds.
Reciprocal checking and prototype updates require
\(O(RMpd_f)\). Exact prototype-neighbour search has a worst-case cost of
\(O(RM^{2}d_f)\), making prototype search the main non-linear term when the
fine partition contains many regions. The PCA decomposition cost depends on
the numerical solver; projection onto a fixed \(d\)-dimensional representation
requires \(O(NBd)\).

The principal memory requirement is

\begin{equation}
O\!\left(NB+E+Nq+Md_f\right),
\label{eq:memory_complexity}
\end{equation}

covering the spectral representation, sparse graph, capped fine-neighbour
graph, and region prototypes. With fixed local graph-neighbourhood caps, graph
construction and sparse diffusion scale approximately linearly with the number
of pixels. The complete implementation does not have linear complexity,
because its exact global fine-neighbour search has quadratic dependence on
\(N\). Exact prototype-neighbour search can introduce an additional non-linear
cost when the fine partition contains many regions.

\subsubsection{Fragmentation and cluster-count controls}

Several diagnostics were used to determine whether high information-based
agreement arose only from producing many regions. First, the final Graph-RMS
partitions were rescored using AMI, homogeneity, completeness, and V-measure.
The difference between NMI and AMI measured the effect of chance correction,
while the homogeneity--completeness balance showed whether the output consisted
mainly of pure but fragmented semantic classes.

Second, a post hoc oracle-\(K\) diagnostic examined how the metrics changed
when the Graph-RMS partition was consolidated to the reference class count.
This experiment used \(K\) only after the partition had been frozen
and was not part of the proposed operating protocol.

Third, MiniBatch-KMeans was evaluated with the same nominal number of clusters
as Graph-RMS. The primary control used
\(\widehat{M}_{\mathrm{full}}\), the number of regions in the complete frozen
Graph-RMS partition. A secondary diagnostic used
\(\widehat{M}_{\Omega}\), the number of Graph-RMS regions intersecting the
reference support. Both controls used standardised PCA-20 features and five
fixed seeds. The labelled-support budget was treated only as a diagnostic
because it depends on the geometry of the reference mask, even though it does
not use the class identities.

\subsubsection{Runtime and memory profiling}

End-to-end Graph-RMS runtime and memory use were profiled on all nine scenes.
Each scene was executed three times on an NVIDIA A100 GPU. The recorded runtime
included data preparation, graph construction, affinity calculation, diffusion,
fine-mode extraction, spatial cleanup, prototype consolidation, and endpoint
selection.

Peak process resident-set size, peak allocated GPU memory, and peak reserved
GPU memory were stored separately. The main manuscript reports process memory
and allocated GPU memory because reserved memory includes allocator capacity
that may not be occupied by active tensors.

An empirical relation between runtime and full-scene pixel count was estimated
from the nine profiled scenes. This fitted trend was treated as a descriptive
scaling summary rather than a formal complexity bound because band count,
graph density, diffusion checkpoint, image geometry, and consolidation
behaviour also differed across scenes.

\FloatBarrier
%==============================================================================
% 3. Results
%==============================================================================
\section{Results}
\label{sec:results}

%==============================================================================
% 3.1 Full-scene performance and classical comparison
%==============================================================================
%==============================================================================
% 3.1 Full-scene performance and classical comparison
%==============================================================================
\subsection{Full-scene performance and classical comparison}
\label{subsec:overall_results}

Graph-RMS returned a full-scene partition for each of the nine datasets without
receiving the reference class count. Its nine-scene mean NMI and ARI were
0.7245 and 0.6086, respectively. The corresponding one-to-one Hungarian mean
OA and BA were 0.6358 and 0.5319. NMI ranged from 0.5885 on Indian Pines to
0.8659 on Salinas-A, while ARI ranged from 0.3022 on Indian Pines to 0.9447
on Trento. Eight of these scenes belonged to the development/evaluation set;
the nine-scene aggregates are therefore reported as descriptive summaries and
not as independent evidence of general superiority. Trento provides the
independent holdout evaluation for the primary procedure.

The number of discovered regions varied more widely than the reference class
count (Table~\ref{tab:primary_results}). Salinas-A produced seven full-scene
regions, six of which intersected the reference mask. KSC produced 720 regions
over the complete image, but only 31 intersected labelled pixels. The
difference is especially relevant for KSC and Botswana, where the reference
masks cover only 1.66\% and 0.86\% of the complete scenes.

\begin{table}[!t]
\centering
\small
\caption{Frozen Graph-RMS results on the nine complete hyperspectral scenes.
\(M_{\mathrm{full}}\) is the number of regions in the complete partition and
\(M_{\Omega}\) is the number intersecting the labelled evaluation support.
OA and BA use one-to-one Hungarian alignment. Trento was the independent
holdout for the primary procedure.}
\label{tab:primary_results}

\renewcommand{\arraystretch}{1.05}
\setlength{\tabcolsep}{4.5pt}

\begin{tabular}{@{}lrrrrrr@{}}
\toprule
Dataset &
\(M_{\mathrm{full}}\) &
\(M_{\Omega}\) &
OA &
BA &
NMI &
ARI \\
\midrule
Salinas-A          & 7   & 6  & 0.8057 & 0.7475 & 0.8659 & 0.7495 \\
Indian Pines       & 35  & 33 & 0.4181 & 0.4956 & 0.5885 & 0.3022 \\
KSC                & 720 & 31 & 0.4642 & 0.3755 & 0.6306 & 0.3226 \\
Pavia University   & 207 & 76 & 0.6365 & 0.5009 & 0.6836 & 0.5137 \\
WHU-Hi LongKou     & 110 & 77 & 0.8110 & 0.5497 & 0.8079 & 0.8261 \\
WHU-Hi HongHu      & 46  & 46 & 0.7704 & 0.6198 & 0.7649 & 0.8586 \\
WHU-Hi HanChuan    & 83  & 79 & 0.5130 & 0.3828 & 0.6134 & 0.6365 \\
Botswana           & 38  & 21 & 0.4498 & 0.4381 & 0.7166 & 0.3234 \\
Trento             & 68  & 41 & 0.8533 & 0.6774 & 0.8495 & 0.9447 \\
\midrule
Nine-scene mean    & --  & -- & 0.6358 & 0.5319 & 0.7245 & 0.6086 \\
\bottomrule
\end{tabular}
\end{table}

Table~\ref{tab:classical_complete} gives the complete comparison with the four
classical methods supplied with the reference value of \(K\). Graph-RMS
obtained the highest OA on six scenes and the highest NMI on all nine. The BA
results were more mixed: Graph-RMS led on Indian Pines, HongHu, HanChuan, and
Trento, while the classical methods led on the remaining scenes. Graph-RMS
also had the highest ARI on six datasets.

\begin{table}[!t]
\centering
\scriptsize
\caption{Complete comparison with classical clustering methods supplied with
the reference class count. Values for PCA-KMeans, MiniBatch-KMeans, FCM, and
SLIC-KMeans are five-seed means. The best value within each dataset and metric
is shown in bold. PCA-KM denotes PCA-KMeans and MB-KM denotes
MiniBatch-KMeans.}
\label{tab:classical_complete}

\renewcommand{\arraystretch}{0.92}
\setlength{\tabcolsep}{5pt}

\begin{tabular}{@{}llrrrr@{}}
\toprule
Dataset & Method & OA & BA & NMI & ARI \\
\midrule

\multirow{5}{*}{Salinas-A}
& Graph-RMS & 0.8057 & 0.7475 & \textbf{0.8659} & 0.7495 \\
& PCA-KM    & 0.7704 & 0.7543 & 0.7695 & 0.6912 \\
& MB-KM     & 0.8001 & 0.8370 & 0.7862 & 0.7076 \\
& FCM       & \textbf{0.8410} & 0.8683 & 0.7828 & 0.7298 \\
& SLIC-KM   & 0.8267 & \textbf{0.8798} & 0.8523 & \textbf{0.7708} \\
\addlinespace[2pt]

\multirow{5}{*}{Indian Pines}
& Graph-RMS & \textbf{0.4181} & \textbf{0.4956} & \textbf{0.5885} & \textbf{0.3022} \\
& PCA-KM    & 0.3745 & 0.4076 & 0.4470 & 0.2258 \\
& MB-KM     & 0.3608 & 0.3649 & 0.4314 & 0.2129 \\
& FCM       & 0.3182 & 0.3315 & 0.4060 & 0.1826 \\
& SLIC-KM   & 0.4044 & 0.3954 & 0.4814 & 0.2307 \\
\addlinespace[2pt]

\multirow{5}{*}{KSC}
& Graph-RMS & 0.4642 & 0.3755 & \textbf{0.6306} & 0.3226 \\
& PCA-KM    & 0.4150 & 0.2544 & 0.4743 & 0.2544 \\
& MB-KM     & 0.4695 & 0.3649 & 0.5335 & 0.3677 \\
& FCM       & \textbf{0.5244} & \textbf{0.4094} & 0.5752 & \textbf{0.4404} \\
& SLIC-KM   & 0.4139 & 0.3251 & 0.5482 & 0.3226 \\
\addlinespace[2pt]

\multirow{5}{*}{Pavia University}
& Graph-RMS & \textbf{0.6365} & 0.5009 & \textbf{0.6836} & \textbf{0.5137} \\
& PCA-KM    & 0.5655 & \textbf{0.5679} & 0.5566 & 0.3560 \\
& MB-KM     & 0.5530 & 0.5518 & 0.5460 & 0.3387 \\
& FCM       & 0.5184 & 0.5446 & 0.5348 & 0.3066 \\
& SLIC-KM   & 0.4808 & 0.4772 & 0.4803 & 0.2837 \\
\addlinespace[2pt]

\multirow{5}{*}{WHU-Hi LongKou}
& Graph-RMS & \textbf{0.8110} & 0.5497 & \textbf{0.8079} & \textbf{0.8261} \\
& PCA-KM    & 0.6987 & 0.5086 & 0.6528 & 0.6931 \\
& MB-KM     & 0.6961 & 0.5454 & 0.6512 & 0.6912 \\
& FCM       & 0.6095 & 0.5101 & 0.6238 & 0.5930 \\
& SLIC-KM   & 0.7479 & \textbf{0.5807} & 0.7133 & 0.7446 \\
\addlinespace[2pt]

\multirow{5}{*}{WHU-Hi HongHu}
& Graph-RMS & \textbf{0.7704} & \textbf{0.6198} & \textbf{0.7649} & \textbf{0.8586} \\
& PCA-KM    & 0.3414 & 0.3449 & 0.4443 & 0.2124 \\
& MB-KM     & 0.3260 & 0.3329 & 0.4378 & 0.2084 \\
& FCM       & 0.2782 & 0.3157 & 0.4207 & 0.1660 \\
& SLIC-KM   & 0.3523 & 0.3888 & 0.5058 & 0.2377 \\
\addlinespace[2pt]

\multirow{5}{*}{WHU-Hi HanChuan}
& Graph-RMS & \textbf{0.5130} & \textbf{0.3828} & \textbf{0.6134} & \textbf{0.6365} \\
& PCA-KM    & 0.4167 & 0.2979 & 0.4105 & 0.3349 \\
& MB-KM     & 0.4276 & 0.3114 & 0.4361 & 0.3873 \\
& FCM       & 0.4122 & 0.3306 & 0.4366 & 0.3927 \\
& SLIC-KM   & 0.4384 & 0.3107 & 0.4308 & 0.3545 \\
\addlinespace[2pt]

\multirow{5}{*}{Botswana}
& Graph-RMS & 0.4498 & 0.4381 & \textbf{0.7166} & 0.3234 \\
& PCA-KM    & 0.5748 & 0.5740 & 0.6052 & 0.3963 \\
& MB-KM     & 0.5097 & 0.5235 & 0.5860 & 0.3666 \\
& FCM       & 0.4745 & 0.4790 & 0.5490 & 0.3321 \\
& SLIC-KM   & \textbf{0.6206} & \textbf{0.6292} & 0.7092 & \textbf{0.5043} \\
\addlinespace[2pt]

\multirow{5}{*}{Trento}
& Graph-RMS & \textbf{0.8533} & \textbf{0.6774} & \textbf{0.8495} & \textbf{0.9447} \\
& PCA-KM    & 0.6451 & 0.5697 & 0.4943 & 0.3636 \\
& MB-KM     & 0.5991 & 0.5289 & 0.4777 & 0.3375 \\
& FCM       & 0.5362 & 0.4972 & 0.4652 & 0.3133 \\
& SLIC-KM   & 0.7816 & 0.5867 & 0.6334 & 0.6830 \\
\midrule

\multirow{5}{*}{Mean}
& Graph-RMS & \textbf{0.6358} & \textbf{0.5319} & \textbf{0.7245} & \textbf{0.6086} \\
& PCA-KM    & 0.5336 & 0.4755 & 0.5394 & 0.3920 \\
& MB-KM     & 0.5269 & 0.4845 & 0.5429 & 0.4020 \\
& FCM       & 0.5014 & 0.4763 & 0.5327 & 0.3840 \\
& SLIC-KM   & 0.5629 & 0.5082 & 0.5950 & 0.4591 \\
\bottomrule
\end{tabular}
\end{table}

The average-rank analysis followed the same pattern
(Table~\ref{tab:classical_statistics}). Graph-RMS had the lowest average rank
for each metric and a rank of 1.000 for NMI. The Friedman \(p\)-values were
0.000205 for NMI and 0.0460 for ARI, compared with 0.0639 for OA and 0.5394
for BA. Given the nine-scene sample, these tests were used to summarise the
cross-scene ranking pattern rather than as the primary evidence for method
performance.

\begin{table}[!t]
\centering
\footnotesize
\caption{Supporting rank and statistical summary for the classical
five-method comparison. Panel A reports average ranks and Friedman tests;
lower rank is better.
Panel B reports Holm-adjusted \(p\)-values for paired comparisons between
Graph-RMS and each classical baseline.}
\label{tab:classical_statistics}

\renewcommand{\arraystretch}{1.02}
\setlength{\tabcolsep}{4.2pt}

\begin{tabular}{@{}lrrrrrrr@{}}
\toprule
\multicolumn{8}{@{}l}{\textit{Panel A: Average ranks and Friedman tests}} \\
\addlinespace[2pt]
Metric &
Graph-RMS &
PCA-KM &
MB-KM &
FCM &
SLIC-KM &
\(\chi^{2}_{F}\) &
\(p\)-value \\
\midrule
OA  & \textbf{1.8889} & 3.2222 & 3.4444 & 3.8889 & 2.5556 & 8.8889  & 0.0639 \\
BA  & \textbf{2.4444} & 3.3333 & 3.2222 & 3.4444 & 2.5556 & 3.1111  & 0.5394 \\
NMI & \textbf{1.0000} & 3.5556 & 3.6667 & 4.1111 & 2.6667 & 21.9556 & 0.000205 \\
ARI & \textbf{1.7778} & 3.4444 & 3.4444 & 3.7778 & 2.5556 & 9.6889  & 0.0460 \\
\midrule

\multicolumn{8}{@{}l}{\textit{Panel B: Holm-adjusted paired \(p\)-values}} \\
\addlinespace[2pt]
Metric &
PCA-KM &
MB-KM &
FCM &
SLIC-KM &
\multicolumn{3}{c}{} \\
\midrule
OA  & 0.5195 & 0.3906 & 0.4375 & 0.8203 & \multicolumn{3}{c}{} \\
BA  & 1.0000 & 1.0000 & 1.0000 & 1.0000 & \multicolumn{3}{c}{} \\
NMI & 0.0625 & 0.0625 & 0.0625 & 0.0625 & \multicolumn{3}{c}{} \\
ARI & 0.2344 & 0.3906 & 0.3008 & 0.5195 & \multicolumn{3}{c}{} \\
\bottomrule
\end{tabular}
\end{table}

Graph-RMS ranked first for NMI on all nine scenes and had the lowest average
rank for each of the four metrics. None of the 16
Graph-RMS-versus-baseline pairwise comparisons remained significant after Holm
correction. The cross-scene tests are therefore treated as supporting
summaries of the observed ranking pattern rather than as evidence of pairwise
statistical superiority.

The Trento result is examined separately in
Section~\ref{subsec:trento_results} because that scene was not used during
method development.

%==============================================================================
% 3.2 Comparison with class-count-free methods
%==============================================================================
\subsection{Comparison with class-count-free methods}
\label{subsec:no_k_results}

The class-count-free comparison produced two different outcomes. Leiden
remained competitive when applied to the same weighted graph as Graph-RMS,
whereas the fixed-default HDBSCAN configuration produced lower
agreement on most scenes. Table~\ref{tab:no_k_complete} reports the absolute
metric values so that both patterns can be inspected directly.

\begin{table}[!t]
\centering
\scriptsize
\caption{Complete comparison of the class-count-free methods. Leiden values
are five-seed means. Graph-RMS and HDBSCAN each contributed one frozen result
under the declared protocol. The best value within each dataset and metric is
shown in bold.}
\label{tab:no_k_complete}

\renewcommand{\arraystretch}{0.94}
\setlength{\tabcolsep}{5pt}

\begin{tabular}{@{}llrrrr@{}}
\toprule
Dataset & Method & OA & BA & NMI & ARI \\
\midrule

\multirow{3}{*}{Salinas-A}
& Graph-RMS & 0.8057 & 0.7475 & 0.8659 & 0.7495 \\
& Leiden    & \textbf{0.8740} & \textbf{0.9010} & \textbf{0.9177} & \textbf{0.8738} \\
& HDBSCAN   & 0.4215 & 0.4420 & 0.5231 & 0.2345 \\
\addlinespace[2pt]

\multirow{3}{*}{Indian Pines}
& Graph-RMS & \textbf{0.4181} & 0.4956 & 0.5885 & \textbf{0.3022} \\
& Leiden    & 0.3866 & \textbf{0.6102} & \textbf{0.6960} & 0.2928 \\
& HDBSCAN   & 0.2410 & 0.0727 & 0.0061 & 0.0012 \\
\addlinespace[2pt]

\multirow{3}{*}{KSC}
& Graph-RMS & 0.4642 & 0.3755 & 0.6306 & 0.3226 \\
& Leiden    & \textbf{0.5426} & \textbf{0.6079} & \textbf{0.7709} & \textbf{0.5091} \\
& HDBSCAN   & 0.1357 & 0.0813 & 0.2856 & -0.0040 \\
\addlinespace[2pt]

\multirow{3}{*}{Pavia University}
& Graph-RMS & \textbf{0.6365} & \textbf{0.5009} & \textbf{0.6836} & \textbf{0.5137} \\
& Leiden    & 0.4365 & 0.4268 & 0.6615 & 0.3710 \\
& HDBSCAN   & 0.4360 & 0.1112 & 0.0002 & 0.0001 \\
\addlinespace[2pt]

\multirow{3}{*}{WHU-Hi LongKou}
& Graph-RMS & \textbf{0.8110} & 0.5497 & \textbf{0.8079} & \textbf{0.8261} \\
& Leiden    & 0.3626 & \textbf{0.5921} & 0.6698 & 0.2931 \\
& HDBSCAN   & 0.3283 & 0.1125 & 0.0032 & 0.0010 \\
\addlinespace[2pt]

\multirow{3}{*}{WHU-Hi HongHu}
& Graph-RMS & \textbf{0.7704} & \textbf{0.6198} & \textbf{0.7649} & \textbf{0.8586} \\
& Leiden    & 0.3651 & 0.5298 & 0.6153 & 0.1771 \\
& HDBSCAN   & 0.4240 & 0.0495 & 0.0064 & 0.0026 \\
\addlinespace[2pt]

\multirow{3}{*}{WHU-Hi HanChuan}
& Graph-RMS & \textbf{0.5130} & 0.3828 & 0.6134 & \textbf{0.6365} \\
& Leiden    & 0.3928 & \textbf{0.4654} & \textbf{0.6595} & 0.3854 \\
& HDBSCAN   & 0.2938 & 0.0759 & 0.0099 & 0.0014 \\
\addlinespace[2pt]

\multirow{3}{*}{Botswana}
& Graph-RMS & 0.4498 & 0.4381 & 0.7166 & 0.3234 \\
& Leiden    & \textbf{0.4575} & \textbf{0.4466} & \textbf{0.7658} & \textbf{0.4960} \\
& HDBSCAN   & 0.1087 & 0.1008 & 0.0407 & 0.0040 \\
\addlinespace[2pt]

\multirow{3}{*}{Trento}
& Graph-RMS & \textbf{0.8533} & \textbf{0.6774} & \textbf{0.8495} & \textbf{0.9447} \\
& Leiden    & 0.6321 & 0.5526 & 0.7617 & 0.7179 \\
& HDBSCAN   & 0.4015 & 0.2736 & 0.2024 & 0.0977 \\
\midrule

\multirow{3}{*}{Mean}
& Graph-RMS & \textbf{0.6358} & 0.5319 & \textbf{0.7245} & \textbf{0.6086} \\
& Leiden    & 0.4944 & \textbf{0.5703} & 0.7242 & 0.4574 \\
& HDBSCAN   & 0.3100 & 0.1466 & 0.1197 & 0.0376 \\
\bottomrule
\end{tabular}
\end{table}

Graph-RMS and Leiden had nearly identical mean NMI values, differing by less
than 0.001. Graph-RMS had higher mean OA and ARI, whereas Leiden had higher
mean BA. The scene-level pattern was more informative than the aggregate mean.
Leiden led on all four metrics for Salinas-A, KSC, and Botswana, while
Graph-RMS led on all four for Pavia University, HongHu, and Trento.

The magnitude of the differences also varied substantially by scene. Relative
to Leiden, Graph-RMS increased OA by 0.4484 and ARI by 0.5330 on LongKou,
OA by 0.4053 and ARI by 0.6815 on HongHu, and OA by 0.2212 and ARI by
0.2268 on Trento. Other scenes favoured Leiden. The post-graph stages
therefore changed the resulting partition materially on several datasets
rather than simply reproducing direct community detection.

None of the paired Graph-RMS--Leiden differences was statistically significant
(Table~\ref{tab:no_k_statistics}). The comparison does not establish a
uniform accuracy advantage for either method. Instead, it shows that the
shared graph contains substantial scene structure, while the subsequent
Graph-RMS diffusion, fine-mode formation, consolidation, and endpoint
selection can produce a different operating point with large gains on some
scenes and losses on others.

\begin{table}[!t]
\centering
\footnotesize
\caption{Scene-level summary of the class-count-free comparisons. A win denotes
a scene on which Graph-RMS obtained the higher metric value. The reported
\(p\)-values are from exact two-sided paired Wilcoxon tests over the nine
scenes.}
\label{tab:no_k_statistics}

\renewcommand{\arraystretch}{1.05}
\setlength{\tabcolsep}{5pt}

\begin{tabular}{@{}llrrrr@{}}
\toprule
Comparator & Quantity & OA & BA & NMI & ARI \\
\midrule
\multirow{2}{*}{Leiden}
& Graph-RMS wins & 6/9 & 3/9 & 4/9 & 6/9 \\
& \(p\)-value    & 0.0977 & 0.4258 & 1.0000 & 0.2031 \\
\addlinespace[2pt]
\multirow{2}{*}{HDBSCAN}
& Graph-RMS wins & 9/9 & 9/9 & 9/9 & 9/9 \\
& \(p\)-value    & 0.003906 & 0.003906 & 0.003906 & 0.003906 \\
\bottomrule
\end{tabular}
\end{table}

Graph-RMS exceeded fixed-default HDBSCAN on all 36 matched scene--metric
comparisons. Under that configuration, HDBSCAN produced only three full-scene
labels on Indian Pines and four on Pavia University. It also assigned 70.0\%
of KSC and 92.4\% of Trento to noise. These results explain the weak agreement
on several scenes, but they apply only to the declared fixed-default HDBSCAN
protocol and not to every possible scene-specific configuration.
%==============================================================================
% 3.3 Holdout performance on Trento
%==============================================================================
%==============================================================================
% 3.3 Holdout performance on Trento
%==============================================================================
\subsection{Holdout performance on Trento}
\label{subsec:trento_results}

Trento was processed after the operating procedure and endpoint rule had been
fixed on the eight development scenes. Its reference map and semantic class
count were not used to revise the method or its held-out configuration. The
holdout run used \(T=100\), \(r_f=0.20\), and \(\tau=0.75\). The partition-based selector
accepted the size-aware branch with \(\beta=0.5\).

The diffused representation produced 213 fine regions. Spatial refinement and
reciprocal prototype consolidation reduced this partition to 68 full-scene
regions, of which 41 intersected the labelled evaluation support. The complete
holdout configuration and selected output are summarised in
Table~\ref{tab:trento_configuration}.

\begin{table}[!t]
\centering
\footnotesize
\caption{Frozen configuration and selected Graph-RMS output for the
independent Trento holdout in the primary experiment.}
\label{tab:trento_configuration}

\renewcommand{\arraystretch}{1.05}
\setlength{\tabcolsep}{5pt}

\begin{tabular}{@{}ccccccc@{}}
\toprule
\(T\) &
\(r_f\) &
\(\tau\) &
\(\beta\) &
Fine regions &
\(M_{\mathrm{full}}\) &
\(M_{\Omega}\) \\
\midrule
100 & 0.20 & 0.75 & 0.5 & 213 & 68 & 41 \\
\bottomrule
\end{tabular}
\end{table}

Graph-RMS attained an NMI of 0.8495 and ARI of 0.9447. Under one-to-one
Hungarian alignment, the corresponding OA and BA were 0.8533 and 0.6774
(Table~\ref{tab:trento_comparison}). Spatial block bootstrapping gave
comparatively narrow intervals for all four measures, with the strongest
agreement observed for ARI.

\begin{table}[!t]
\centering
\footnotesize
\caption{Performance on the independent Trento holdout for the primary
Graph-RMS procedure. Graph-RMS intervals are
95\% spatial block-bootstrap intervals from 300 replicates with a block size
of 16 pixels. SLIC-KMeans was the strongest classical method on Trento.
Leiden values are five-seed means, and HDBSCAN used the declared fixed-default
configuration.}
\label{tab:trento_comparison}

\renewcommand{\arraystretch}{1.08}
\setlength{\tabcolsep}{5pt}

\begin{tabular}{@{}lrrrr@{}}
\toprule
Method or summary & OA & BA & NMI & ARI \\
\midrule
\textbf{Graph-RMS}
& \textbf{0.8533}
& \textbf{0.6774}
& \textbf{0.8495}
& \textbf{0.9447} \\

95\% bootstrap interval
& [0.8219, 0.8949]
& [0.6452, 0.7383]
& [0.8288, 0.8826]
& [0.9135, 0.9664] \\

SLIC-KMeans
& 0.7816
& 0.5867
& 0.6334
& 0.6830 \\

Leiden
& 0.6321
& 0.5526
& 0.7617
& 0.7179 \\

HDBSCAN
& 0.4015
& 0.2736
& 0.2024
& 0.0977 \\
\bottomrule
\end{tabular}
\end{table}

Relative to SLIC-KMeans, Graph-RMS improved OA by 0.0717, BA by 0.0907,
NMI by 0.2161, and ARI by 0.2617. It also exceeded both class-count-free
comparators on the holdout. Trento had not been used to revise the primary procedure before this
evaluation.

The information-based diagnostics in Table~\ref{tab:trento_diagnostics}
clarify the structure of the selected partition. Homogeneity and purity were
high, indicating that individual regions were usually dominated by one
reference category. Completeness was lower because several reference
categories were represented by more than one discovered region.

\begin{table}[!t]
\centering
\footnotesize
\caption{Structural and information-based diagnostics for the selected Trento
partition. The final two quantities describe how many discovered regions were
needed to account for 90\% of a reference class.}
\label{tab:trento_diagnostics}

\renewcommand{\arraystretch}{1.07}
\setlength{\tabcolsep}{7pt}

\begin{tabular}{@{}lr@{}}
\toprule
Diagnostic & Value \\
\midrule
Adjusted mutual information & 0.8492 \\
Homogeneity & 0.9636 \\
Completeness & 0.7596 \\
Many-to-one purity & 0.9877 \\
Mean regions covering 90\% of a class & 4.67 \\
Maximum regions covering 90\% of a class & 11 \\
\bottomrule
\end{tabular}
\end{table}

The contrast between high homogeneity and lower completeness is consistent
with a fine structural partition rather than a direct reconstruction of the
six-class reference legend. The relation between structural regions and
semantic classes is examined across all nine scenes in the next subsection.


%==============================================================================
% 3.4 Structural granularity and semantic correspondence
%==============================================================================
\subsection{Structural granularity and semantic correspondence}
\label{subsec:fragmentation_results}

Graph-RMS did not constrain the number of discovered regions to match the
number of reference classes. The selected full-scene granularity ranged from
seven regions on Salinas-A to 720 on KSC. This variation reflects both scene
structure and the fact that Graph-RMS processed unlabelled areas outside the
reference masks.

Figure~\ref{fig:representative_maps} shows three representative outcomes. On
KSC, the reference labels are sparse and spatially discontinuous, so many
full-scene regions do not appear within the evaluation mask. HongHu contains
larger labelled areas and shows closer spatial correspondence, although some
fields are divided into finer regions. Trento preserves the main reference
structures while retaining subdivisions within several semantic categories.

\begin{figure}[!t]
\centering
\includegraphics[
    width=0.68\textwidth,
    height=0.36\textheight,
    keepaspectratio
]{figures/fig2_representative_cluster_maps.png}
\caption{Representative Graph-RMS partitions and reference maps for KSC,
WHU-Hi HongHu, and Trento. Region identifiers were aligned with reference
classes only for display and evaluation. Black pixels have no reference label,
while grey pixels belong to regions that remained unmatched under the
one-to-one Hungarian assignment. This alignment was not used during graph
construction, diffusion, consolidation, or endpoint selection.}
\label{fig:representative_maps}
\end{figure}

Table~\ref{tab:region_granularity} distinguishes the number of regions in the
complete image from the number intersecting the labelled support. The
difference is largest for KSC: its 720 full-scene regions correspond to 55.38
regions per reference class, but only 31 regions intersect the reference mask,
giving a labelled-support ratio of 2.38. Botswana shows a smaller version of
the same pattern.

\begin{table}[!t]
\centering
\footnotesize
\caption{Structural granularity of the selected Graph-RMS partitions.
\(K\) is the reference class count, \(M_{\mathrm{full}}\) is the complete
full-scene region count, and \(M_{\Omega}\) is the number of regions
intersecting the labelled support. The ratios are descriptive and were not
used during endpoint selection.}
\label{tab:region_granularity}

\renewcommand{\arraystretch}{1.04}
\setlength{\tabcolsep}{4.5pt}

\begin{tabular}{@{}lrrrrr@{}}
\toprule
Dataset &
\(K\) &
\(M_{\mathrm{full}}\) &
\(M_{\Omega}\) &
\(M_{\mathrm{full}}/K\) &
\(M_{\Omega}/K\) \\
\midrule
Salinas-A        & 6  & 7   & 6  & 1.17  & 1.00 \\
Indian Pines     & 16 & 35  & 33 & 2.19  & 2.06 \\
KSC              & 13 & 720 & 31 & 55.38 & 2.38 \\
Pavia University & 9  & 207 & 76 & 23.00 & 8.44 \\
WHU-Hi LongKou   & 9  & 110 & 77 & 12.22 & 8.56 \\
WHU-Hi HongHu    & 22 & 46  & 46 & 2.09  & 2.09 \\
WHU-Hi HanChuan  & 16 & 83  & 79 & 5.19  & 4.94 \\
Botswana         & 14 & 38  & 21 & 2.71  & 1.50 \\
Trento           & 6  & 68  & 41 & 11.33 & 6.83 \\
\bottomrule
\end{tabular}
\end{table}

Pavia University, LongKou, and Trento retained the finest structure within
their labelled areas, with \(M_{\Omega}/K\) values of 8.44, 8.56, and 6.83,
respectively. Salinas-A was the closest to a one-region-per-class result, with
six label-intersecting regions for six reference classes.

Figure~\ref{fig:fragmentation} relates these region counts to the difference
between many-to-one purity and one-to-one OA. The latter penalises a partition
when several predicted regions correspond to the same reference class, even
when each individual region is internally consistent.

\begin{figure}[!t]
\centering
\includegraphics[
    width=0.76\textwidth,
    height=0.34\textheight,
    keepaspectratio
]{figures/fig3_fragmentation.png}
\caption{Structural granularity and semantic fragmentation across the nine
scenes. The upper panel compares full-scene and labelled-support region counts
with the reference class count on a logarithmic scale. The lower panel shows
the difference between many-to-one purity and one-to-one Hungarian OA. A
larger difference indicates that several class-consistent regions correspond
to the same reference category.}
\label{fig:fragmentation}
\end{figure}

Many-to-one purity was consistently higher than one-to-one OA. On Trento,
purity was 0.9877, compared with an OA of 0.8533. Indian Pines and HanChuan
showed the largest purity--OA differences, at 0.2655 and 0.2543,
respectively. Their mean numbers of regions needed to cover 90\% of a
reference class were 3.13 and 6.38. Pavia University required a mean of 4.56
regions, with a maximum of 19 for one reference category. These results show
that the one-to-one accuracy penalty arose mainly from semantic classes being
divided into several regions.

Two controls examined whether the strong information-based scores could be
explained only by the number of discovered regions
(Table~\ref{tab:granularity_controls}). First, MiniBatch-KMeans was rerun using
the Graph-RMS full-scene region count as its requested number of clusters.
Graph-RMS retained higher OA and NMI on all nine scenes, higher BA on seven,
and higher ARI on eight. Second, the frozen Trento partition was consolidated
post hoc to the six reference classes. This reduced all four primary metrics.

\begin{table}[!t]
\centering
\footnotesize
\caption{Controls for region-count and semantic-fragmentation effects.
Panel A compares the selected Trento partition with a post hoc six-region
version and a count-matched MiniBatch-KMeans result. Dashes indicate metrics
not included in the stored count-matched summary. Panel B gives the number of
scenes on which Graph-RMS exceeded count-matched MiniBatch-KMeans.}
\label{tab:granularity_controls}

\renewcommand{\arraystretch}{1.06}
\setlength{\tabcolsep}{5pt}

\begin{tabular}{@{}lrrrr@{}}
\toprule
\multicolumn{5}{@{}l}{\textit{Panel A: Trento controls}} \\
\addlinespace[2pt]
Partition or control & OA & BA & NMI & ARI \\
\midrule
Selected Graph-RMS
& 0.8533 & 0.6774 & 0.8495 & 0.9447 \\

Graph-RMS consolidated to \(K=6\)
& 0.7916 & 0.5691 & 0.7611 & 0.7267 \\

Count-matched MiniBatch-KMeans
& -- & -- & 0.4108 & 0.0974 \\
\midrule

\multicolumn{5}{@{}l}{\textit{Panel B: Graph-RMS wins over count-matched
MiniBatch-KMeans}} \\
\addlinespace[2pt]
Comparison & OA & BA & NMI & ARI \\
\midrule
Scenes won & 9/9 & 7/9 & 9/9 & 8/9 \\
\bottomrule
\end{tabular}
\end{table}

The decrease after forcing Trento to six regions indicates that the selected
partition retained spectral--spatial distinctions that were lost under direct
compression to the semantic class count. Similarly, the count-matched
MiniBatch-KMeans control shows that a large region count alone does not account
for the observed NMI and ARI.

The Graph-RMS region count should therefore be read as the granularity of the
discovered spectral--spatial structure. It is not an estimate of the number of
thematic classes in the reference map.
%==============================================================================
% 3.5 Component ablation
%==============================================================================
%==============================================================================
% 3.5 Component ablation
%==============================================================================
\subsection{Component ablation}
\label{subsec:ablation_results}

The Pavia University ablation examined how each stage affected the scale and
reference agreement of the final partition. Each variant changed one part of
the frozen Graph-RMS procedure while retaining the remaining settings. The
complete method produced 207 full-scene regions, with an OA of 0.6365, BA of
0.5009, NMI of 0.6836, and ARI of 0.5137
(Table~\ref{tab:ablation}).

\begin{table}[!t]
\centering
\footnotesize
\caption{Component ablation on Pavia University. The region count refers to
the complete scene. Each variant changed one component of the frozen
Graph-RMS procedure. The highest value in each metric column is shown in
bold.}
\label{tab:ablation}

\renewcommand{\arraystretch}{1.05}
\setlength{\tabcolsep}{4.2pt}

\begin{tabularx}{\textwidth}{@{}Xrrrrr@{}}
\toprule
Configuration & Regions & OA & BA & NMI & ARI \\
\midrule
Complete Graph-RMS
& 207
& \textbf{0.6365}
& 0.5009
& \textbf{0.6836}
& 0.5137 \\

Without mutual-neighbour gate
& 33
& 0.6125
& \textbf{0.6424}
& 0.6374
& 0.3495 \\

Without spatial affinity
& 1
& 0.4360
& 0.1111
& 0.0000
& 0.0000 \\

Without diffusion (\(T=0\))
& 7\,336
& 0.0129
& 0.0266
& 0.2943
& 0.0022 \\

Without spatial cleanup
& 2\,214
& 0.6125
& 0.4418
& 0.6495
& \textbf{0.5156} \\

Without fine-mode grouping
& 189\,178
& 0.0022
& 0.0050
& 0.2863
& 0.0001 \\

Without prototype consolidation
& 692
& 0.6194
& 0.4525
& 0.6602
& 0.5117 \\

Without size weighting (\(\beta=0\))
& 213
& 0.5548
& 0.5477
& 0.5157
& 0.2413 \\
\bottomrule
\end{tabularx}
\end{table}

The graph affinity and diffusion stages controlled opposite forms of
degeneracy. Removing the spatial affinity connected the scene too broadly and
left a single region, reducing NMI and ARI to zero. Removing diffusion
preserved too much pixel-level variation: the partition expanded to 7\,336
regions, with an OA of 0.0129 and an ARI of 0.0022. The later consolidation
stages could not recover a useful partition from this fragmented
representation.

Fine-mode grouping had an even larger effect on granularity. Without it,
189\,178 regions remained, and all four reference-based measures were low.
Spatial cleanup reduced the number of isolated assignments from 2\,214 to 207
after subsequent processing. Its removal caused only a small change in ARI
but reduced OA, BA, and NMI. The complete method and the no-cleanup variant
had ARI values of 0.5137 and 0.5156, respectively, showing that the cleanup
stage primarily improved spatial regularity and one-to-one semantic agreement
rather than every partition metric.

Prototype consolidation reduced the cleaned fine partition from 692 to 207
regions. Omitting consolidation produced moderate declines in OA, BA, and NMI,
while ARI remained close to that of the complete method. The result suggests
that many fine regions already followed useful structural boundaries, but
consolidation produced a more manageable partition without discarding most of
that agreement.

The size adjustment affected which prototypes merged. Setting \(\beta=0\)
retained a similar number of regions and increased BA from 0.5009 to 0.5477,
but reduced OA from 0.6365 to 0.5548, NMI from 0.6836 to 0.5157, and ARI
from 0.5137 to 0.2413. Removing the mutual-neighbour gate also increased BA,
but reduced the other three metrics and compressed the scene to 33 regions.
Neither modification was uniformly inferior, but both shifted the partition
towards a different balance between class recall and structural agreement.

On Pavia University, the ablation exposed two principal failure modes.
Removing major graph, diffusion, or fine-mode operations led either to collapse
or extensive fragmentation, while cleanup and prototype consolidation altered
the scale of the retained regions. The complete configuration did not maximise
every metric, but it avoided the degenerate outputs observed after removal of
the major stages and retained the highest OA and NMI among the tested
variants. The experiment used one diagnostic scene, so these component effects
should be interpreted as illustrative rather than exhaustive across all
datasets.


%==============================================================================
% 3.6 Parameter sensitivity and endpoint-selector behaviour
%==============================================================================
\subsection{Parameter sensitivity and endpoint-selector behaviour}
\label{subsec:sensitivity_results}

The sensitivity analysis comprised 51 runs on Salinas-A, Indian Pines, and
Pavia University. The spatial-window radius, PCA dimension used for neighbour
ranking, diffusion checkpoint, and fine-mode radius were varied one at a time,
with the remaining settings held at their frozen values. Figure
\ref{fig:sensitivity} reports the corresponding OA and NMI values.

\begin{figure}[!t]
\centering
\includegraphics[
    width=0.90\textwidth,
    height=0.38\textheight,
    keepaspectratio
]{figures/fig4_sensitivity.png}
\caption{Sensitivity of OA and NMI to the PCA search dimension, spatial-window
radius, diffusion checkpoint, and fine-mode radius on Salinas-A, Indian Pines,
and Pavia University. Each parameter was varied separately while the remaining
settings were held fixed. Single-region outputs are retained in the figure
rather than removed from the tested range.}
\label{fig:sensitivity}
\end{figure}

The principal responses are summarised in
Table~\ref{tab:sensitivity_summary}. The PCA search dimension produced the
smallest changes on Salinas-A and Pavia University. Salinas-A retained seven
regions across 10--40 components, with OA between 0.8057 and 0.8089 and NMI
between 0.8649 and 0.8772. Pavia University retained 207--221 regions, with
NMI between 0.6637 and 0.6836. Indian Pines was more sensitive: using ten
components increased OA to 0.5678 and ARI to 0.4668, compared with frozen
values of 0.4181 and 0.3022 at 20 components.

\begin{table}[!t]
\centering
\footnotesize
\caption{Summary of the main parameter-sensitivity responses. The table
reports the most informative changes from the complete 51-run analysis;
the full OA and NMI trajectories are shown in
Fig.~\ref{fig:sensitivity}.}
\label{tab:sensitivity_summary}

\renewcommand{\arraystretch}{1.06}
\setlength{\tabcolsep}{4pt}

\begin{tabularx}{\textwidth}{@{}p{0.18\textwidth}
p{0.16\textwidth}
p{0.17\textwidth}
X@{}}
\toprule
Parameter & Scene & Tested range & Main response \\
\midrule

PCA search dimension
& Salinas-A
& 10--40
& Seven regions throughout; OA \(0.8057\)--\(0.8089\) and
NMI \(0.8649\)--\(0.8772\). \\

PCA search dimension
& Indian Pines
& 10--40
& Ten components gave OA \(0.5678\) and ARI \(0.4668\), compared with
the frozen OA \(0.4181\) and ARI \(0.3022\). \\

PCA search dimension
& Pavia University
& 10--40
& Region count \(207\)--\(221\); NMI \(0.6637\)--\(0.6836\). \\

Spatial-window radius
& Salinas-A
& 3, 5, 7, 9
& Radii 5, 7, and 9 all produced seven regions with similar agreement. \\

Spatial-window radius
& Indian Pines
& 3, 5, 7, 9
& Final region count varied from 30 to 87 without a single-region collapse. \\

Spatial-window radius
& Pavia University
& 3, 5, 7, 9
& Radius 5 produced the highest OA and NMI; radius 9 produced one region. \\

Diffusion checkpoint
& Salinas-A
& 25, 50, 100, 200
& Agreement improved from \(T=25\) to \(T=100\) and changed little at
\(T=200\). \\

Diffusion checkpoint
& Indian Pines
& 12, 25, 50, 100
& OA was highest at \(T=25\); the frozen \(T=50\) retained the highest
NMI among the longer checkpoints. \\

Diffusion checkpoint
& Pavia University
& 25, 50, 100, 200
& Checkpoints 25--100 produced 148--207 regions; \(T=200\) produced one
region. \\

Fine-mode radius
& Salinas-A
& 0.075--0.225
& Results were comparatively stable apart from the lower response at
\(r_f=0.1125\). \\

Fine-mode radius
& Indian Pines
& 0.150--0.450
& NMI declined from \(0.6102\) at \(r_f=0.225\) to \(0.4017\) at
\(r_f=0.450\). \\

Fine-mode radius
& Pavia University
& 0.075--0.225
& Increasing \(r_f\) from 0.1875 to 0.225 changed the output from
46 regions to one region. \\
\bottomrule
\end{tabularx}
\end{table}

The spatial window, diffusion checkpoint, and fine-mode radius had stronger
scene-dependent effects than the PCA search dimension. Pavia University was
the most susceptible to abrupt consolidation. A large spatial window, a long
diffusion trajectory, or a sufficiently broad fine-mode radius could connect
structures that remained separate under the frozen setting. Indian Pines
changed more gradually, while Salinas-A remained stable over most of the
tested range.

The endpoint selector was examined separately from the parameter sweep. It
returned a reported partition for all nine unperturbed scenes. The size-aware
branch was accepted on eight scenes, while WHU-Hi LongKou used the conservative
fallback. Relative to a retrospective metric-specific oracle over the stored
candidate endpoints, the selector chose the highest-OA candidate on seven
scenes. Pavia University, Botswana, and Trento matched the oracle endpoint for
all four primary metrics.

\begin{table}[!t]
\centering
\footnotesize
\caption{Endpoint-selector audit. Panel A reports the mean difference between
the selected endpoint and the retrospective best candidate for each metric.
Panel B summarises the main rule-ablation outcomes. Reference metrics were
used only after the partition-based endpoint had been selected.}
\label{tab:selector_audit}

\renewcommand{\arraystretch}{1.06}
\setlength{\tabcolsep}{4.5pt}

\begin{tabularx}{\textwidth}{@{}lrrrr@{}}
\toprule
\multicolumn{5}{@{}l}{\textit{Panel A: Retrospective oracle gap}} \\
\addlinespace[2pt]
Quantity & OA & BA & NMI & ARI \\
\midrule
Mean oracle gap
& 0.0042
& 0.0206
& 0.0093
& 0.0058 \\
\midrule
\multicolumn{5}{@{}l}{\textit{Panel B: Selector-rule ablations}} \\
\addlinespace[2pt]
\multicolumn{1}{@{}l}{Rule change} &
\multicolumn{4}{l}{Observed consequence} \\
\midrule

\multicolumn{1}{@{}l}{Remove largest-region constraint}
&
\multicolumn{4}{p{0.70\textwidth}@{}}{
Indian Pines decreased from 35 to 10 regions; HongHu decreased from 46 to
six regions. HongHu BA, NMI, and ARI fell to 0.2153, 0.6076, and 0.6073.
} \\

\multicolumn{1}{@{}l}{Use stability alone}
&
\multicolumn{4}{p{0.70\textwidth}@{}}{
LongKou decreased to two regions; HongHu and Botswana collapsed to one
region. Mean NMI fell from 0.7245 to 0.4665.
} \\

\multicolumn{1}{@{}l}{Remove entropy condition}
&
\multicolumn{4}{p{0.70\textwidth}@{}}{
No selected endpoint changed in the nine nominal runs.
} \\

\multicolumn{1}{@{}l}{Remove compression condition}
&
\multicolumn{4}{p{0.70\textwidth}@{}}{
No selected endpoint changed in the nine nominal runs.
} \\

\multicolumn{1}{@{}l}{Remove persistence condition}
&
\multicolumn{4}{p{0.70\textwidth}@{}}{
Only HongHu changed, from 46 to 49 regions, with negligible metric
differences.
} \\
\bottomrule
\end{tabularx}
\end{table}

The mean oracle gaps were 0.0042 for OA, 0.0206 for BA, 0.0093 for NMI,
and 0.0058 for ARI. These values indicate that the selected endpoints were
usually close to the best stored candidate, although the larger BA gap shows
that the partition-based criteria were not equivalent to optimising class-balanced
recall.

The rule ablation identified the largest-region constraint as the main defence
against excessive consolidation. Removing it reduced Indian Pines from 35 to
10 regions and HongHu from 46 to six. On HongHu, BA fell from 0.6198 to
0.2153, NMI from 0.7649 to 0.6076, and ARI from 0.8586 to 0.6073.
A stability-only selector was still less reliable: it reduced LongKou to two
regions and produced single-region outputs on HongHu and Botswana. Its mean
NMI was 0.4665, compared with 0.7245 for the complete selector.

Removing the entropy or compression condition did not alter any nominal
endpoint, while removing persistence changed only HongHu from 46 to 49
regions with negligible metric differences. These conditions were inactive in
most nominal runs but remained part of the rule because they describe
different failure modes and become relevant when the consolidation path
changes under perturbation.

The sensitivity runs did not identify one operating scale that behaved
similarly on all three scenes. PCA search dimension had the smallest effect,
whereas spatial reach, diffusion depth, and fine-mode radius could change the
partition abruptly. The endpoint guards rejected several collapsed solutions,
but the primary procedure still required scene-specific scale settings. This
dependence motivated the automatic-selection experiment described next.
%==============================================================================
% 3.7 Perturbation response and cross-device repeatability
%==============================================================================
%==============================================================================
% 3.7 Perturbation response and cross-device repeatability
%==============================================================================
\subsection{Automatic operating-point selection}
\label{subsec:auto_selector_results}

The locked automatic rule returned an admissible operating point for all eight
development scenes. Seven scenes selected the conservative \(\beta=0\)
branch, while KSC selected the size-aware branch. No development scene
required abstention.

Across the eight development scenes, automatic selection gave mean OA, BA,
NMI, and ARI values of 0.5744, 0.5384, 0.7239, and 0.5813, respectively.
The corresponding means from the primary scene-specific configurations were
0.6086, 0.5137, 0.7089, and 0.5666. Automatic selection therefore changed
the mean metrics by \(-0.0341\) for OA, \(+0.0246\) for BA,
\(+0.0150\) for NMI, and \(+0.0147\) for ARI.

\begin{table}[!t]
\centering
\footnotesize
\caption{Mean performance across the eight development scenes for the primary
scene-specific configurations and the automatic-selection experiment.
Reference metrics were calculated only after the automatic operating point
for each scene had been selected.}
\label{tab:auto_selector_summary}

\renewcommand{\arraystretch}{1.06}
\setlength{\tabcolsep}{7pt}

\begin{tabular}{@{}lrrrr@{}}
\toprule
Procedure & OA & BA & NMI & ARI \\
\midrule
Primary scene-specific
& 0.6086 & 0.5137 & 0.7089 & 0.5666 \\
Automatic selection
& 0.5744 & 0.5384 & 0.7239 & 0.5813 \\
Difference
& -0.0341 & +0.0246 & +0.0150 & +0.0147 \\
\bottomrule
\end{tabular}
\end{table}

The retrospective Trento transfer test selected \(T=100\),
\(\gamma=1.25\), \(r_f=0.1360\), \(\tau=0.75\), and \(\beta=0.5\),
producing 67 full-scene regions. After the selected partition had been frozen,
reference-based evaluation gave an OA of 0.6148, BA of 0.5122, NMI of 0.7656,
and ARI of 0.7164.

The automatic experiment therefore showed that the principal operating
settings could be selected from image-derived partition statistics without
using target-scene labels or \(K\). The procedure retained comparable
information-based agreement across the development scenes, with higher mean
BA, NMI, and ARI and lower mean OA than the primary scene-specific
configurations. The Trento automatic result is reported as a retrospective
transfer test rather than as a second independent holdout.


\subsection{Perturbation response and cross-device repeatability}
\label{subsec:robustness_results}

Perturbation response was examined on Salinas-A, Indian Pines, and Pavia
University using two independently seeded runs for 1\% additive noise and
5\% random pixel replacement. Salinas-A and Indian Pines retained their frozen
consolidation thresholds. For Pavia University, the complete partition-based
endpoint rule was rerun because the original threshold was not assumed to
remain admissible after perturbation.

Table~\ref{tab:perturbation_results} reports the complete semantic response.
Salinas-A changed little under additive noise. Its two NMI values were 0.8668
and 0.8668, compared with 0.8659 for the frozen partition, while ARI remained
near 0.750. Pixel replacement produced eight and seven regions and retained
similar or slightly higher NMI and ARI values.

\begin{table}[!t]
\centering
\scriptsize
\caption{Perturbation response on Salinas-A, Indian Pines, and Pavia
University. Salinas-A and Indian Pines used their frozen consolidation
thresholds. Pavia University used partition-based endpoint re-selection. The
unperturbed row is included as the reference for each scene.}
\label{tab:perturbation_results}

\renewcommand{\arraystretch}{0.98}
\setlength{\tabcolsep}{3.5pt}

\begin{tabular}{@{}llclrrrrr@{}}
\toprule
Dataset &
Condition &
Run &
Selector status &
Regions &
OA &
BA &
NMI &
ARI \\
\midrule

\multirow{5}{*}{Salinas-A}
& Unperturbed & -- & Frozen result & 7
& 0.8057 & 0.7475 & 0.8659 & 0.7495 \\
& Noise 1\% & 1 & Frozen threshold & 7
& 0.8068 & 0.7484 & 0.8668 & 0.7506 \\
& Noise 1\% & 2 & Frozen threshold & 7
& 0.8063 & 0.7479 & 0.8668 & 0.7506 \\
& Pixel replacement 5\% & 1 & Frozen threshold & 8
& 0.8087 & 0.7492 & 0.8871 & 0.7719 \\
& Pixel replacement 5\% & 2 & Frozen threshold & 7
& 0.8110 & 0.7516 & 0.8785 & 0.7617 \\
\addlinespace[2pt]

\multirow{5}{*}{Indian Pines}
& Unperturbed & -- & Frozen result & 35
& 0.4181 & 0.4956 & 0.5885 & 0.3022 \\
& Noise 1\% & 1 & Frozen threshold & 36
& 0.4188 & 0.5361 & 0.5899 & 0.2996 \\
& Noise 1\% & 2 & Frozen threshold & 34
& 0.4430 & 0.5244 & 0.5864 & 0.3042 \\
& Pixel replacement 5\% & 1 & Frozen threshold & 41
& 0.3853 & 0.4820 & 0.5795 & 0.2542 \\
& Pixel replacement 5\% & 2 & Frozen threshold & 39
& 0.3249 & 0.4127 & 0.5108 & 0.1349 \\
\addlinespace[2pt]

\multirow{5}{*}{Pavia University}
& Unperturbed & -- & Frozen result & 207
& 0.6365 & 0.5009 & 0.6836 & 0.5137 \\
& Noise 1\% & 1 & Selected & 102
& 0.6275 & 0.5214 & 0.6813 & 0.4933 \\
& Noise 1\% & 2 & Selected & 147
& 0.6390 & 0.5123 & 0.6873 & 0.5133 \\
& Pixel replacement 5\% & 1 & Abstained & 700
& 0.4760 & 0.3098 & 0.4943 & 0.2399 \\
& Pixel replacement 5\% & 2 & Abstained & 697
& 0.4723 & 0.3098 & 0.4964 & 0.2372 \\
\bottomrule
\end{tabular}
\end{table}

Indian Pines remained close to its frozen result under additive noise, although
the two runs differed more in OA and BA than Salinas-A. Pixel replacement had
a clearer effect. NMI decreased to 0.5795 and 0.5108, while ARI decreased to
0.2542 and 0.1349. The second replacement run therefore produced the largest
semantic change among the fixed-threshold trials.

For Pavia University, the selector accepted both noise-perturbed sequences.
The selected outputs contained 102 and 147 regions, and their NMI values of
0.6813 and 0.6873 were close to the unperturbed value of 0.6836. Their
full-scene ARI values against the frozen partition were 0.9628 and 0.9813;
the corresponding labelled-support values were 0.9636 and 0.9974
(Table~\ref{tab:pavia_partition_agreement}).

\begin{table}[!t]
\centering
\footnotesize
\caption{Partition-level agreement for the adaptive Pavia University
perturbation runs. ARI was calculated between each perturbed partition and the
frozen unperturbed partition.}
\label{tab:pavia_partition_agreement}

\renewcommand{\arraystretch}{1.05}
\setlength{\tabcolsep}{5pt}

\begin{tabular}{@{}lclrrr@{}}
\toprule
Perturbation &
Run &
Selector status &
Regions &
ARI, all pixels &
ARI, labelled pixels \\
\midrule
Noise 1\% & 1 & Selected & 102 & 0.9628 & 0.9636 \\
Noise 1\% & 2 & Selected & 147 & 0.9813 & 0.9974 \\
Pixel replacement 5\% & 1 & Abstained & 700 & 0.4855 & 0.5563 \\
Pixel replacement 5\% & 2 & Abstained & 697 & 0.4851 & 0.5782 \\
\bottomrule
\end{tabular}
\end{table}

Neither pixel-replacement sequence contained a persistent, non-degenerate
Pavia endpoint. The selector therefore abstained and retained the
unconsolidated diagnostic outputs rather than presenting them as accepted
Graph-RMS partitions. These outputs contained 700 and 697 regions and had
full-scene ARI values of approximately 0.485 against the frozen partition.
The abstention records a loss of endpoint support rather than hiding the
perturbation response behind a forced result.

\begin{figure}[!t]
\centering
\includegraphics[
    width=0.86\textwidth,
    height=0.32\textheight,
    keepaspectratio
]{figures/repeatability_and_perturbation_matched_style.png}
\caption{Perturbation response and cross-device repeatability. The left panel
shows ARI between frozen CPU partitions and new GPU partitions over all pixels
and over the labelled support. The right panel shows agreement between the
frozen Pavia University partition and the adaptively re-selected perturbation
outputs. The selector accepted both additive-noise trials and abstained for
both pixel-replacement trials.}
\label{fig:robustness_repeatability}
\end{figure}

Cross-device repeatability was examined by comparing the frozen CPU partitions
with new GPU executions on KSC, WHU-Hi HanChuan, and Botswana
(Table~\ref{tab:cross_device_repeatability}). KSC and Botswana retained
all-pixel ARI values above 0.99. HanChuan showed the largest change, with its
region count increasing from 83 to 102 and its all-pixel ARI decreasing to
0.9114. Agreement remained higher on its labelled support, where ARI was
0.9578.

\begin{table}[!t]
\centering
\footnotesize
\caption{Cross-device partition repeatability. Frozen CPU partitions were
compared with new GPU executions. Region-count differences do not imply that
the principal spatial partition changed to the same extent.}
\label{tab:cross_device_repeatability}

\renewcommand{\arraystretch}{1.06}
\setlength{\tabcolsep}{5pt}

\begin{tabular}{@{}lrrrrr@{}}
\toprule
Dataset &
CPU regions &
GPU regions &
Difference &
ARI, all pixels &
ARI, labelled pixels \\
\midrule
KSC
& 720 & 686 & \(-34\) & 0.9960 & 0.9996 \\

WHU-Hi HanChuan
& 83 & 102 & \(+19\) & 0.9114 & 0.9578 \\

Botswana
& 38 & 43 & \(+5\) & 0.9921 & 0.9854 \\
\bottomrule
\end{tabular}
\end{table}

The cross-device results show strong partition-level repeatability, but not
bitwise identity. Most differences were associated with the splitting or
merging of small regions. Graph-RMS should therefore be expected to preserve
the principal scene structure across the tested computational paths without
guaranteeing an identical region count.


%==============================================================================
% 3.8 Runtime, memory use, and empirical scaling
%==============================================================================
\subsection{Runtime, memory use, and empirical scaling}
\label{subsec:runtime_results}

Runtime and memory use were measured over three complete executions of each
scene on an NVIDIA A100-SXM4 GPU with 80~GB of memory. The timing boundary
included scene standardisation, graph construction, affinity calculation,
diffusion, fine-mode formation, spatial cleanup, prototype consolidation, and
endpoint selection.

Table~\ref{tab:runtime_memory_complete} reports the full profiling results.
Mean runtime ranged from \(1.21\)~s on Salinas-A to \(42.34\)~s on
WHU-Hi HongHu. Trento required \(4.54\)~s, while the three largest runtime
values occurred on HongHu, HanChuan, and Botswana.

\begin{table}[!t]
\centering
\footnotesize
\caption{End-to-end Graph-RMS runtime and memory use on the nine complete
scenes. Runtime values are means and standard deviations over three executions
on an NVIDIA A100-SXM4 GPU. Process memory is peak resident-set size; GPU
memory is peak allocated memory.}
\label{tab:runtime_memory_complete}

\renewcommand{\arraystretch}{1.05}
\setlength{\tabcolsep}{5pt}

\begin{tabular}{@{}lrrr@{}}
\toprule
Dataset &
Runtime (s) &
Peak process memory (MB) &
Peak allocated GPU memory (MB) \\
\midrule
Salinas-A
& \(1.21 \pm 0.56\)
& 1\,538
& 253 \\

Indian Pines
& \(1.60 \pm 0.05\)
& 2\,005
& 716 \\

KSC
& \(23.34 \pm 3.20\)
& 4\,007
& 10\,518 \\

Pavia University
& \(11.34 \pm 0.06\)
& 3\,021
& 6\,944 \\

WHU-Hi LongKou
& \(21.61 \pm 2.89\)
& 3\,522
& 7\,364 \\

WHU-Hi HongHu
& \(42.34 \pm 4.38\)
& 5\,084
& 14\,932 \\

WHU-Hi HanChuan
& \(37.28 \pm 3.70\)
& 4\,593
& 12\,335 \\

Botswana
& \(24.74 \pm 2.93\)
& 4\,498
& 12\,639 \\

Trento
& \(4.54 \pm 0.58\)
& 2\,433
& 3\,342 \\
\bottomrule
\end{tabular}
\end{table}

\begin{figure}[!t]
\centering
\includegraphics[
    width=0.86\textwidth,
    height=0.32\textheight,
    keepaspectratio
]{figures/fig6_runtime_memory.png}
\caption{End-to-end runtime and peak allocated GPU memory for the nine
complete scenes. Runtime values are three-run means. The memory values show
active allocation rather than the larger amount reserved by the GPU memory
allocator.}
\label{fig:runtime_memory}
\end{figure}

Peak allocated GPU memory ranged from 253~MB for Salinas-A to 14\,932~MB
for HongHu. KSC, HanChuan, and Botswana each required more than 10~GB, while
Trento required 3.34~GB. Peak process memory was lower than allocated GPU
memory on the larger scenes and ranged from approximately 1.54 to 5.08~GB.

Resource use was not determined by pixel count alone. Botswana contained fewer
spectral bands than the three WHU-Hi scenes but required 12.64~GB of allocated
GPU memory. LongKou required 7.36~GB despite having more pixels than Pavia
University. Differences in spectral dimensionality, retained graph edges,
diffusion depth, image geometry, fine-region count, and consolidation path all
contributed to the observed profile.

Across the nine scenes, the fitted relation between mean runtime and pixel
count was

\begin{equation}
t(N)\propto N^{0.891}.
\label{eq:empirical_runtime_scaling}
\end{equation}

The sublinear exponent describes the measured scene set and should not be read
as a theoretical complexity result. The formal analysis in
Section~\ref{subsec:additional_analysis} shows that sparse graph construction
and diffusion are approximately linear in \(N\) when the local graph
neighbourhood caps are fixed. The current exact global fine-neighbour search,
however, has quadratic dependence on \(N\), and prototype-neighbour search can
also become expensive when the fine partition contains many regions. The
empirical fit therefore describes only the measured runtime trend over the
nine tested scenes and does not imply near-linear complexity for the complete
implementation.

All nine primary scenes completed in less than 45~s on the A100, making runtime
less restrictive than memory use for the tested scenes. Peak allocated GPU
memory nevertheless reached approximately 15~GB on HongHu. This requirement
can exceed the capacity of commodity GPUs and becomes important when extending
the method to larger swaths. Lower-memory diffusion, streamed affinity
construction, and approximate prototype-neighbour search would therefore be
needed for more memory-constrained or substantially larger deployments.

\FloatBarrier
%==============================================================================
% 7. Discussion
%==============================================================================
%==============================================================================
% 4. Discussion
%==============================================================================
%==============================================================================
% 4. Discussion
%==============================================================================
\section{Discussion}
\label{sec:discussion}

Graph-RMS recovered spectral--spatial structure that corresponded with the
reference maps without receiving target-scene labels or the reference class
count during processing. Its most consistent result was NMI: the method ranked
first on all nine scenes in the five-method classical protocol. OA and BA were more variable because they used
one-to-one region-to-class assignment. When several coherent regions belonged
to the same semantic category, only one could be matched to that category and
the remaining regions were counted as errors. The difference between these
metric families is therefore part of the result rather than a contradiction
between them.

The independent Trento holdout supports this interpretation. Graph-RMS reached
an NMI of 0.8495 and ARI of 0.9447 after the procedure had been fixed on the
development scenes; the corresponding one-to-one Hungarian OA was 0.8533. Its homogeneity of 0.9636 shows that most discovered
regions were dominated by one reference category, while the lower completeness
of 0.7596 indicates that several categories were represented by multiple
regions. Forcing the frozen partition to six groups reduced all four primary
metrics. The additional regions therefore retained useful structure rather
than arising only from arbitrary label proliferation.

The comparison with Leiden clarifies the contribution of the stages that
follow graph construction. The nearly identical mean NMI values show that the
weighted graph already captures a substantial part of the recoverable scene
structure, and the present evidence does not support claiming that Graph-RMS
is uniformly superior to direct community detection. The two procedures,
however, do not simply produce equivalent outcomes. Graph-RMS gave markedly
higher OA and ARI on LongKou, HongHu, and Trento, whereas Leiden was stronger
on Salinas-A, KSC, and Botswana. These scene-dependent differences indicate
that diffusion and subsequent region formation alter how the graph structure
is expressed rather than merely adding computation to the same partition.

The additional stages also provide an explicit fine-to-coarse operating
procedure. Diffusion changes the representation before region formation,
fine modes expose local structure, reciprocal prototype consolidation reduces
that structure, and the endpoint rule rejects collapsed or unsupported
candidates. The Pavia University ablation shows that removing these stages can
produce either collapse or severe fragmentation, while the perturbation
experiments show that the endpoint rule can abstain when an admissible
partition is not supported. These properties do not constitute a general
accuracy advantage over Leiden, but they define the operating behaviour that
distinguishes Graph-RMS from direct community detection
\citep{traag2019}.

Fixed-default HDBSCAN behaved differently. Graph-RMS exceeded it in all 36
scene--metric comparisons, while HDBSCAN assigned large fractions of KSC and
Trento to noise and returned very few groups on several scenes. This result is
specific to the declared fixed-default protocol. HDBSCAN was included as a
density-based reference and was not exhaustively optimised for each scene, so
the comparison should not be interpreted as general superiority over
density-based clustering \citep{campello2013}.

The Pavia University ablation shows why Graph-RMS must be considered as an
integrated procedure. Removing spatial affinity produced a single region,
whereas removing diffusion or fine-mode grouping produced thousands of small
regions. Spatial cleanup reduced isolated assignments, and prototype
consolidation converted the detailed fine partition into a more usable region
set. The size-aware dispersion term improved OA, NMI, and ARI relative to the
conservative endpoint, although BA moved in the opposite direction. No single
stage improved every measure; together, the stages controlled collapse,
fragmentation, and unsupported merging.

The sensitivity analysis identifies scale selection as the main source of
variation in the primary procedure. PCA dimension had a comparatively small
effect on Salinas-A and Pavia University, whereas spatial reach, diffusion
checkpoint, and fine-mode radius could alter the partition abruptly. The
largest-region constraint prevented several stable but degenerate endpoints,
and stability alone was insufficient because a collapsed partition could
remain stable over adjacent consolidation thresholds.

The automatic-selection experiment addressed this dependence directly. After
a single global rule had been calibrated on the eight development scenes,
selection for an individual scene required neither its reference labels nor
its semantic class count. Mean BA, NMI, and ARI increased by 0.0246, 0.0150,
and 0.0147 relative to the primary scene-specific configurations, while mean
OA decreased by 0.0341. The result indicates that manually specified
scene-specific scales can be replaced by a common partition-based selection
procedure without reducing the mean information-based agreement on the
development scenes. The lower OA is consistent with the finer partitions
returned by the automatic procedure and the one-to-one matching penalty used
for that metric.

The automatic Trento run provides a retrospective transfer test of the same
rule. It selected an operating point without using Trento labels during
selection and obtained NMI \(=0.7656\) and ARI \(=0.7164\), compared with
0.8495 and 0.9447 for the pre-specified primary Trento configuration. The
corresponding OA decreased from 0.8533 to 0.6148. The difference confirms
that scale selection remains consequential: automatic selection removes the
need to prescribe the principal target-scene scales, but it does not reproduce
the performance of the manually specified primary configuration on Trento.
Since the primary Trento result was already known when the automatic rule was
designed, this run is interpreted as a retrospective transfer test rather than
an additional independent holdout. Prospective evaluation on further unseen
acquisitions is needed to establish how broadly the automatic rule transfers
across sensors and scene types.

Perturbation behaviour also depended on scene structure. Salinas-A remained
close to its frozen partition under both tested perturbations, while Indian
Pines was more affected by pixel replacement. On Pavia University, endpoint
re-selection recovered partitions close to the frozen result under additive
noise and abstained under the stronger pixel-replacement trials. This
abstention is useful because it distinguishes a supported endpoint from a
partition produced merely because the pipeline was forced to return one.
Cross-device ARI remained above 0.91 on the three tested scenes, although
small-region counts changed.

Graph-RMS is most directly suited to exploratory analysis before a reliable
semantic legend is available. The discovered regions may support
candidate-region inspection or subsequent annotation, although annotation
efficiency was not measured here. A further limitation is that a final
Graph-RMS region label need not correspond to a single spatially connected
component. Fine-mode links and prototype consolidation operate in the
diffusion-feature space, so spatially separated components can receive the
same final region label even though spatial information is incorporated during
graph construction and cleanup. Computational scalability remains a practical
constraint. Peak allocated GPU memory reached approximately 15~GB on HongHu.
The current exact global fine-neighbour search has quadratic dependence on the
number of pixels, while exact prototype-neighbour search has quadratic
worst-case dependence on the number of current prototypes. The 189\,178-region output in
the Pavia ablation occurred after fine-mode grouping was deliberately removed
and therefore does not represent normal operation, but it illustrates why
controlling fine-partition granularity is computationally important. The
automatic-selection experiment reduces dependence on manually specified scene
scales, while prospective transfer across sensors and acquisition conditions
remains to be established. Lower-memory graph operations, approximate
prototype search, and broader transfer testing are therefore the main
practical directions for extending the present method.
% 5. Conclusion
%==============================================================================
\section{Conclusion}
\label{sec:conclusion}

Graph-RMS provides a training-free procedure for discovering coherent
spectral--spatial regions without requiring labelled training samples or a
predefined semantic class count during scene processing. In the primary
nine-scene evaluation, it achieved the highest NMI on every scene in the
five-method classical comparison. On the independent Trento holdout, NMI and
ARI were 0.8495 and 0.9447, while the corresponding one-to-one Hungarian OA
was 0.8533. Graph-RMS exceeded fixed-default HDBSCAN under the declared
protocol, whereas Leiden remained a competitive same-graph alternative and
showed that graph construction itself accounts for a substantial part of the
recoverable structure.

A separate automatic-selection experiment showed that the principal
scene-dependent operating settings could be selected without target-scene
labels or \(K\) after a single global rule had been calibrated on the
development scenes. Across those scenes, automatic selection increased mean
BA, NMI, and ARI while reducing mean OA. The same frozen rule returned an
operating point in a retrospective Trento transfer test, although its
reference-based agreement was lower than that of the pre-specified primary
configuration. Graph-RMS is therefore best viewed as a method for discovering
spectral--spatial structure during exploratory analysis rather than for
estimating the number of semantic classes.
%==============================================================================
% Declarations
%==============================================================================


\section*{Author Contributions}

Simhaa T.\,T.: Methodology, Software, Validation, Formal analysis,
Investigation, Visualisation, Writing -- original draft.
Steven Yves Le Moan: Conceptualisation, Methodology, Supervision,
Writing -- review and editing.
Shihabudheen K.\,V.: Conceptualisation, Methodology, Supervision,
Writing -- review and editing.
\section*{Funding}

This research received no specific grant from any funding agency in the
public, commercial, or not-for-profit sectors.

\section*{Disclosure statement}

The authors report there are no competing interests to declare.


\section*{Data availability statement}

The hyperspectral datasets used in this study are publicly available from the
original repositories cited in the manuscript. The Graph-RMS source code,
derived partitions, experimental configurations, evaluation summaries, and
numerical results are available from
\url{https://github.com/Simhaatt/Graph-RMS} and are archived for release
\texttt{v1.0.0} at \url{https://doi.org/10.5281/zenodo.21876448}. The
repository does not redistribute the original hyperspectral datasets.

\section*{Ethics statement}

Not applicable. The study used public remote-sensing benchmark datasets and
did not involve human participants, personal data, or animals.

\section*{Declaration of Generative AI and AI-assisted technologies}

During preparation of this manuscript and its reproducibility package, the
authors used OpenAI ChatGPT and OpenAI Codex (GPT-5 family, accessed
July--August 2026) to assist with code development and debugging, manuscript
organisation, documentation, and language refinement. The tools were not used
to generate experimental observations or to make autonomous scientific
decisions. All source code, experimental outputs, numerical values,
references, figures, and manuscript text were reviewed and verified by the
authors. The authors take full responsibility for the accuracy, originality,
and integrity of the submitted work.
%==============================================================================
% Supplementary material
%==============================================================================

%==============================================================================
% References
%==============================================================================
\bibliographystyle{plainnat}
\bibliography{references}

\end{document}
